% =====================================================================
% Supporting Information for
% "Uncertainty Quantification for Aquatic Ecotoxicity Prediction"
%
% Included as the final section of main.tex via % =====================================================================
% Supporting Information for
% "Uncertainty Quantification for Aquatic Ecotoxicity Prediction"
%
% Included as the final section of main.tex via % =====================================================================
% Supporting Information for
% "Uncertainty Quantification for Aquatic Ecotoxicity Prediction"
%
% Included as the final section of main.tex via % =====================================================================
% Supporting Information for
% "Uncertainty Quantification for Aquatic Ecotoxicity Prediction"
%
% Included as the final section of main.tex via \input{supporting_information}.
% Relies on the preamble of main.tex (amsmath, booktabs, etc. are already
% loaded there; amssymb is added there for \Real below). Uses the shared
% references.bib already cited by the rest of the paper.
% =====================================================================

% ---- convenience macros ----
\newcommand{\yhat}{\hat{y}}
\newcommand{\Norm}{\mathcal{N}}
\newcommand{\Gam}{\mathrm{Gamma}}
\newcommand{\Real}{\mathbb{R}}
\newcommand{\E}{\mathbb{E}}
\newcommand{\Var}{\mathrm{Var}}
\newcommand{\vv}{\mathbf{v}}
\newcommand{\xx}{\mathbf{x}}

\section{Supporting Information}
\label{sec:si}

This Supporting Information gives the full mathematical specification of the
Hierarchical Bayesian Factorization Machine (BFM) summarised in the main text.
It is written for a statistical or machine-learning audience and documents the
model exactly as implemented: the predictor, the complete hierarchical
generative model, the Gibbs sampler and every conditional posterior, the
empirical-Bayes initialisation, the hyperparameter settings, and the
posterior-predictive quantities used for calibration, Species Sensitivity
Distributions (SSDs), and hazardous concentrations. The predictive core is the
factorization machine (FM) of \citet{rendle2010fm} cast in the Bayesian form of
\citet{rendle2011bayesian}; Section~\ref{sec:diff} states precisely how our
model differs from theirs. Symbols follow the main text and are collected in
Table~\ref{tab:notation}.

% =====================================================================
\subsection{Notation}
\label{sec:notation}

\begin{table}[htbp]
\centering
\caption{Notation used throughout the Supporting Information.}
\label{tab:notation}
\small
\begin{tabular}{ll}
\toprule
\textbf{Symbol} & \textbf{Meaning} \\
\midrule
$n$ & number of observations (rows) \\
$p$ & number of features (design-matrix columns) \\
$k$ & number of latent factors per feature (factorization rank) \\
$C$ & number of chemicals (noise groups) \\
$\xx_i \in \Real^{p}$ & feature vector of observation $i$ \\
$x_{ij}$ & $j$-th entry of $\xx_i$ \\
$y_i$ & centred log-transformed LC50 of observation $i$ \\
$c_i \in \{1,\dots,C\}$ & chemical (group) index of observation $i$ \\
$w_0$ & global bias \\
$w_j$ & linear weight of feature $j$ \\
$v_{jf}$ & $f$-th latent factor of feature $j$; $\vv_j=(v_{j1},\dots,v_{jk})$ \\
$\alpha_c$ & noise precision of chemical $c$ (aleatoric term) \\
$\mu_w,\lambda_w$ & mean and precision of the shared prior on the weights $w_j$ \\
$\mu_v,\lambda_v$ & mean and precision of the shared prior on the factors $v_{jf}$ \\
$a_0,b_0$ & shape and rate of the shared prior on the $\alpha_c$ \\
$\mu_0,\gamma_0$ & mean and precision of the Gaussian mean-hyperpriors \\
$\alpha_\lambda,\beta_\lambda$ & shape and rate of the Gamma prior on $\lambda_w,\lambda_v$ \\
$c_0,d_0$ & shape and rate of the Gamma hyperprior on $b_0$ \\
$\theta$ & all model parameters $\{w_0,\{w_j\},\{v_{jf}\}\}$ \\
$S$ & number of retained (post-burn-in) posterior samples \\
$\yhat^{(s)}$ & prediction from posterior sample $s$ \\
\bottomrule
\end{tabular}
\end{table}

We write $\Norm(\mu,\sigma^2)$ for a Gaussian with mean $\mu$ and variance
$\sigma^2$, and $\Gam(a,b)$ for a Gamma distribution with shape $a$ and rate $b$
(so that its mean is $a/b$). All precisions are inverse variances.

% =====================================================================
\subsection{Data representation and feature encoding}
\label{sec:features}

Each observation is a single laboratory LC50 measurement for a
(chemical, species, duration) combination. The target $y_i$ is the
log-transformed concentration in mg/L, centred by subtracting the global mean
$\bar{y}$ of the training set; the mean is added back to recover log-mg/L units
when reporting predictions. Repeated measurements of the same triplet are kept
as separate rows rather than averaged, so that the per-chemical noise model
(Section~\ref{sec:model}) sees the within-triplet spread directly.

The feature vector $\xx_i$ concatenates five one-hot categorical blocks and two
standardised numerical features:
\begin{itemize}
  \item \textbf{chemical} identity (one-hot over $C$ chemicals),
  \item \textbf{species} identity (one-hot),
  \item \textbf{test duration} (one-hot over $\{24,48,72,96\}$\,h),
  \item \textbf{taxonomic family} and \textbf{taxonomic class} of the species (two one-hot blocks),
  \item \textbf{numerical}: log-transformed molecular weight and the RDKit
        octanol--water partition coefficient (\texttt{clogp}), median-imputed
        and standardised to zero mean and unit variance.
\end{itemize}
Unseen categories at prediction time map to an all-zero block (they contribute
nothing to the sums below), so the model can score any (chemical, species,
duration) triplet whose chemical and species were each observed at least once
during training, but not entirely novel chemicals or species. For the dataset
used here $n=70{,}670$, $C=3{,}295$ chemicals, $1{,}267$ species, and
$p=4{,}879$ (namely $3{,}295+1{,}267+4+283+28+2$).

% =====================================================================
\subsection{The factorization-machine predictor}
\label{sec:fm}

The predictive mean for observation $i$ is the second-order factorization
machine
\begin{equation}
  \yhat_i(\theta) \;=\; w_0 \;+\; \sum_{j=1}^{p} w_j\,x_{ij}
  \;+\; \sum_{j=1}^{p}\sum_{j'>j}^{p} \langle \vv_j,\vv_{j'}\rangle\,x_{ij}x_{ij'},
  \label{eq:fm}
\end{equation}
with $\langle \vv_j,\vv_{j'}\rangle=\sum_{f=1}^{k} v_{jf}v_{j'f}$. The pairwise
interactions are not stored explicitly; they are represented in the low-rank
form of Eq.~\eqref{eq:fm}, which lets the model infer the interaction between a
chemical and a species that were never tested together. Using the standard FM
identity, the interaction sum is evaluated in $O(pk)$ time (in practice $O(k
N_{\mathrm{nz}})$ for a sparse $\xx_i$ with $N_{\mathrm{nz}}$ non-zeros) as
\begin{equation}
  \sum_{j<j'} \langle \vv_j,\vv_{j'}\rangle\,x_{ij}x_{ij'}
  \;=\; \tfrac{1}{2}\sum_{f=1}^{k}\!\left[
      \Big(\sum_{j=1}^{p} v_{jf}x_{ij}\Big)^{\!2}
      - \sum_{j=1}^{p} v_{jf}^{2}x_{ij}^{2}
  \right].
  \label{eq:fm-identity}
\end{equation}
The diagonal subtraction in Eq.~\eqref{eq:fm-identity} removes each feature's
self-interaction. A direct consequence, used repeatedly by the sampler, is that
$\yhat_i$ is an \emph{affine} function of each individual parameter when the
others are held fixed (Section~\ref{sec:sampler}).

% =====================================================================
\subsection{Full hierarchical generative model}
\label{sec:model}

The model is a Bayesian FM in which the single global noise precision of the
standard BFM is replaced by a per-chemical precision $\alpha_c$ under a shared
Gamma prior. The complete generative model is:

\paragraph{Likelihood.} Each observation is its FM mean plus Gaussian noise
whose precision is that of the observation's chemical,
\begin{equation}
  y_i \;\mid\; \xx_i,\theta,\alpha_{c_i}
  \;\sim\; \Norm\!\big(\yhat_i(\theta),\ \alpha_{c_i}^{-1}\big),
  \qquad i=1,\dots,n.
  \label{eq:like}
\end{equation}

\paragraph{Priors on the model parameters.} The global bias has a fixed Gaussian
prior; the weights and factors are drawn from shared populations whose mean and
precision are themselves learned:
\begin{align}
  w_0     &\sim \Norm(\mu_0,\ \gamma_0^{-1}), \\
  w_j     &\sim \Norm(\mu_w,\ \lambda_w^{-1}), & j&=1,\dots,p, \label{eq:prior-w}\\
  v_{jf}  &\sim \Norm(\mu_v,\ \lambda_v^{-1}), & j&=1,\dots,p,\ f=1,\dots,k. \label{eq:prior-v}
\end{align}

\paragraph{Hyperpriors on the populations.} There is one population for all
linear weights and one for all latent factors (as in the base BFM of
\citet{rendle2011bayesian}); each has a Gaussian hyperprior on its mean and a
Gamma hyperprior on its precision:
\begin{align}
  \mu_w &\sim \Norm(\mu_0,\ \gamma_0^{-1}), & \lambda_w &\sim \Gam(\alpha_\lambda,\ \beta_\lambda), \\
  \mu_v &\sim \Norm(\mu_0,\ \gamma_0^{-1}), & \lambda_v &\sim \Gam(\alpha_\lambda,\ \beta_\lambda).
\end{align}

\paragraph{Hierarchical noise model.} This is the methodological contribution.
Each chemical has its own noise precision, drawn from a Gamma prior whose rate is
shared across all chemicals and is itself given a Gamma hyperprior:
\begin{align}
  \alpha_c &\sim \Gam(a_0,\ b_0), & c&=1,\dots,C, \label{eq:prior-alpha}\\
  b_0      &\sim \Gam(c_0,\ d_0). \label{eq:prior-b0}
\end{align}
The shape $a_0$ is held fixed (set by empirical Bayes, Section~\ref{sec:eb});
the rate $b_0$ is resampled every sweep. Because $b_0$ is common to all
chemicals, it ties the per-chemical precisions together: a chemical with many
observations has its $\alpha_c$ set mainly by the scatter of its own residuals,
whereas a sparsely measured chemical is pooled towards the population-typical
noise level implied by $\Gam(a_0,b_0)$. This partial pooling is what keeps the
aleatoric estimate stable for data-poor chemicals.

The joint density factorises as
\begin{equation}
\begin{aligned}
  p\big(\theta,\{\alpha_c\},b_0,\mu_w,\lambda_w,\mu_v,\lambda_v \mid \mathbf{y}\big)
  \;\propto\;
  & \Big[\textstyle\prod_{i=1}^{n} \Norm\!\big(y_i \mid \yhat_i(\theta),\alpha_{c_i}^{-1}\big)\Big]
    \Norm(w_0\mid\mu_0,\gamma_0^{-1}) \\
  &\times \Big[\textstyle\prod_{j=1}^{p}\Norm(w_j\mid\mu_w,\lambda_w^{-1})\Big]
          \Big[\textstyle\prod_{j,f}\Norm(v_{jf}\mid\mu_v,\lambda_v^{-1})\Big] \\
  &\times \Norm(\mu_w\mid\mu_0,\gamma_0^{-1})\,\Gam(\lambda_w\mid\alpha_\lambda,\beta_\lambda)\,
          \Norm(\mu_v\mid\mu_0,\gamma_0^{-1})\,\Gam(\lambda_v\mid\alpha_\lambda,\beta_\lambda) \\
  &\times \Big[\textstyle\prod_{c=1}^{C}\Gam(\alpha_c\mid a_0,b_0)\Big]\,\Gam(b_0\mid c_0,d_0).
\end{aligned}
\label{eq:joint}
\end{equation}
This is the model drawn in plate notation in the main text: the fixed constants
$(\mu_0,\gamma_0,\alpha_\lambda,\beta_\lambda,a_0,c_0,d_0)$ set the hyperpriors;
$w_j$ and $v_{jf}$ share $(\mu_w,\lambda_w)$ and $(\mu_v,\lambda_v)$
respectively; the $\alpha_c$ share the rate $b_0$; and each $y_i$ is drawn from a
Gaussian centred at $\yhat_i$ with precision $\alpha_{c_i}$.

% =====================================================================
\subsection{Relationship to the standard Bayesian FM}
\label{sec:diff}

Our model departs from the BFM of \citet{rendle2011bayesian} in two respects.
\begin{enumerate}
  \item \textbf{Per-chemical noise (the substantive change).} The standard BFM
  assumes a single global precision $\alpha$ shared by all observations, with a
  Gamma prior and a Gamma-conjugate posterior. We replace it by $C$
  chemical-specific precisions $\alpha_c$ (Eq.~\ref{eq:prior-alpha}) under a
  shared, learned rate $b_0$ (Eq.~\ref{eq:prior-b0}). Equivalently, the global
  $\alpha$ in every conditional below becomes the observation-specific
  $\alpha_{c_i}$, and one extra conjugate draw for $b_0$ is added per sweep. This
  turns the noise term into a per-chemical, partially pooled quantity, which is
  the aleatoric component carried into the downstream analysis.

  \item \textbf{Independent Gaussian priors on the population means (a minor
  change).} The standard BFM couples each population mean to its precision
  through a Normal--Gamma prior, $\mu \mid \lambda \sim
  \Norm(\mu_0,(\gamma_0\lambda)^{-1})$, so the posterior precision of $\mu$
  scales with $\lambda$. We instead place an independent
  $\Norm(\mu_0,\gamma_0^{-1})$ prior on each mean (Eqs.~\ref{eq:prior-w}--\ref{eq:prior-v}),
  decoupling $\mu$ from $\lambda$; the mean does not then enter the update for
  $\lambda$, and the $\mu$-posterior precision is $\gamma_0+p\lambda$ rather than
  $(\gamma_0+p)\lambda$ (Section~\ref{sec:sampler}). With $p$ in the thousands
  this has negligible effect on inference, as \citet{rendle2011bayesian} also
  note for the hyperprior constants, but we state it for exactness.
\end{enumerate}
As in the reference, we use single-parameter (unblocked) Gibbs sampling, which
avoids the matrix inversions of blocked schemes and keeps each sweep linear in
$k$.

% =====================================================================
\subsection{Gibbs sampler and conditional posteriors}
\label{sec:sampler}

Posterior inference is by Gibbs sampling: each quantity is drawn in turn from its
full conditional given the current values of the others. All conditionals are
available in closed form because (i) the FM mean is affine in each single
parameter, giving Gaussian conditionals with Gaussian priors, and (ii) the
Gamma priors on the precisions are conjugate to the Gaussian likelihood.

Throughout we cache the running residual
\begin{equation}
  e_i \;=\; y_i - \yhat_i(\theta),
  \label{eq:resid}
\end{equation}
and the per-factor cache $q_{if}=\sum_{j=1}^{p} v_{jf}x_{ij}$, both updated by a
rank-one correction after each parameter is resampled. This is what makes a full
sweep $O(k N_{\mathrm{nz}})$, where $N_{\mathrm{nz}}$ is the number of non-zeros
in the design matrix. For a parameter with current value $\theta_\ell$ and
``slope'' $h_{i\ell}=\partial \yhat_i/\partial\theta_\ell$ (constant in
$\theta_\ell$ by multilinearity), the partial residual that adds back its current
contribution is $r_{i}=e_i+h_{i\ell}\theta_\ell$.

\subsubsection{Global bias $w_0$}
Here $h_{i0}=1$ for all $i$, so with prior $\Norm(\mu_0,\gamma_0^{-1})$,
\begin{equation}
  \sigma_{w_0}^2=\Big(\gamma_0+\sum_{i=1}^{n}\alpha_{c_i}\Big)^{-1},
  \qquad
  \mu_{w_0}=\sigma_{w_0}^2\Big(\gamma_0\mu_0+\sum_{i=1}^{n}\alpha_{c_i}\,(e_i+w_0)\Big),
  \qquad
  w_0 \sim \Norm(\mu_{w_0},\sigma_{w_0}^2).
\end{equation}
Then set $e_i \mathrel{+}= (w_0^{\mathrm{old}}-w_0^{\mathrm{new}})$ for all $i$.

\subsubsection{Weight population $(\mu_w,\lambda_w)$}
Treating the $p$ weights as i.i.d.\ $\Norm(\mu_w,\lambda_w^{-1})$ draws,
\begin{align}
  \mu_w &\sim \Norm\!\big(\mu_{\mu_w},\,\sigma_{\mu_w}^2\big),
  & \sigma_{\mu_w}^2 &=\big(\gamma_0+p\,\lambda_w\big)^{-1},
  & \mu_{\mu_w} &=\sigma_{\mu_w}^2\big(\gamma_0\mu_0+\lambda_w\textstyle\sum_{j=1}^{p}w_j\big), \\
  \lambda_w &\sim \Gam\!\Big(\alpha_\lambda+\tfrac{p}{2},\ \
     \beta_\lambda+\tfrac{1}{2}\textstyle\sum_{j=1}^{p}(w_j-\mu_w)^2\Big).
\end{align}

\subsubsection{Linear weights $w_j$}
For feature $j$, $h_{ij}=x_{ij}$, and only observations with $x_{ij}\neq 0$
contribute. With prior $\Norm(\mu_w,\lambda_w^{-1})$ and partial residual
$r_i=e_i+x_{ij}w_j$,
\begin{equation}
  \sigma_{w_j}^2=\Big(\lambda_w+\sum_{i:\,x_{ij}\neq0}\!\alpha_{c_i}\,x_{ij}^2\Big)^{-1},
  \quad
  \mu_{w_j}=\sigma_{w_j}^2\Big(\lambda_w\mu_w+\sum_{i:\,x_{ij}\neq0}\!\alpha_{c_i}\,x_{ij}\,r_i\Big),
  \quad
  w_j\sim\Norm(\mu_{w_j},\sigma_{w_j}^2).
\end{equation}
Then update $e_i \mathrel{+}= x_{ij}\,(w_j^{\mathrm{old}}-w_j^{\mathrm{new}})$.

\subsubsection{Factor population $(\mu_v,\lambda_v)$}
Treating all $pk$ factor entries as i.i.d.\ $\Norm(\mu_v,\lambda_v^{-1})$,
\begin{align}
  \mu_v &\sim \Norm\!\big(\mu_{\mu_v},\,\sigma_{\mu_v}^2\big),
  & \sigma_{\mu_v}^2 &=\big(\gamma_0+p\,k\,\lambda_v\big)^{-1},
  & \mu_{\mu_v} &=\sigma_{\mu_v}^2\big(\gamma_0\mu_0+\lambda_v\textstyle\sum_{j,f}v_{jf}\big), \\
  \lambda_v &\sim \Gam\!\Big(\alpha_\lambda+\tfrac{p k}{2},\ \
     \beta_\lambda+\tfrac{1}{2}\textstyle\sum_{j,f}(v_{jf}-\mu_v)^2\Big).
\end{align}

\subsubsection{Latent factors $v_{jf}$}
By the multilinearity of Eq.~\eqref{eq:fm-identity}, $\yhat_i$ is affine in
$v_{jf}$ with slope
\begin{equation}
  h_{ijf}=x_{ij}\big(q_{if}-x_{ij}v_{jf}\big),
  \qquad q_{if}=\sum_{j'=1}^{p} v_{j'f}x_{ij'},
\end{equation}
which does not depend on $v_{jf}$ (the $x_{ij}v_{jf}$ subtraction removes it).
With prior $\Norm(\mu_v,\lambda_v^{-1})$ and partial residual
$r_i=e_i+h_{ijf}v_{jf}$,
\begin{equation}
  \sigma_{v_{jf}}^2=\Big(\lambda_v+\sum_{i:\,x_{ij}\neq0}\!\alpha_{c_i}\,h_{ijf}^2\Big)^{-1},
  \quad
  \mu_{v_{jf}}=\sigma_{v_{jf}}^2\Big(\lambda_v\mu_v+\sum_{i:\,x_{ij}\neq0}\!\alpha_{c_i}\,h_{ijf}\,r_i\Big),
  \quad
  v_{jf}\sim\Norm(\mu_{v_{jf}},\sigma_{v_{jf}}^2).
\end{equation}
Then update both caches: $e_i \mathrel{+}= h_{ijf}\,(v_{jf}^{\mathrm{old}}-v_{jf}^{\mathrm{new}})$
and $q_{if}\mathrel{+}= x_{ij}\,(v_{jf}^{\mathrm{new}}-v_{jf}^{\mathrm{old}})$.

\subsubsection{Per-chemical noise precisions $\alpha_c$}
Let $N_c=|\{i:c_i=c\}|$ and $\mathrm{SSE}_c=\sum_{i:\,c_i=c} e_i^2$ be the
observation count and residual sum of squares of chemical $c$. Conjugacy of the
Gamma prior (Eq.~\ref{eq:prior-alpha}) to the Gaussian likelihood
(Eq.~\ref{eq:like}) gives
\begin{equation}
  \alpha_c \;\sim\; \Gam\!\Big(a_0+\tfrac{N_c}{2},\ \ b_0+\tfrac{1}{2}\,\mathrm{SSE}_c\Big),
  \qquad c=1,\dots,C.
  \label{eq:post-alpha}
\end{equation}
All $C$ precisions are drawn in one vectorised step using per-group sums of
$e_i^2$.

\subsubsection{Shared noise rate $b_0$}
The rate $b_0$ is conjugate to the $C$ Gamma-distributed precisions with common
shape $a_0$: combining $\prod_c \Gam(\alpha_c\mid a_0,b_0)\propto
b_0^{C a_0}e^{-b_0\sum_c\alpha_c}$ with the prior $\Gam(c_0,d_0)$ yields
\begin{equation}
  b_0 \;\sim\; \Gam\!\Big(c_0+C\,a_0,\ \ d_0+\textstyle\sum_{c=1}^{C}\alpha_c\Big).
  \label{eq:post-b0}
\end{equation}

\subsubsection{One sweep and sample collection}
A single Gibbs sweep applies, in order: $w_0$; $(\mu_w,\lambda_w)$ then every
$w_j$; $(\mu_v,\lambda_v)$ then every $v_{jf}$; the $\alpha_c$
(Eq.~\ref{eq:post-alpha}); and $b_0$ (Eq.~\ref{eq:post-b0}). The chain is run
for $n_{\mathrm{iter}}$ sweeps from an arbitrary initialisation
($w_0=0$, $w=\mathbf{0}$, $v_{jf}\sim\Norm(0,0.1^2)$, $\alpha_c=1$); the first
$n_{\mathrm{burn}}$ sweeps are discarded as burn-in and the remaining
$S=n_{\mathrm{iter}}-n_{\mathrm{burn}}$ sweeps are retained. Each retained sample
$s$ stores $\{w_0^{(s)},w^{(s)},v^{(s)},\{\alpha_c^{(s)}\},b_0^{(s)}\}$ and
corresponds to one complete, internally consistent parameterisation of the
model.

% =====================================================================
\subsection{Empirical-Bayes initialisation of the prior shape $a_0$}
\label{sec:eb}

The shared shape $a_0$ controls how strongly sparsely tested chemicals are
pooled towards the population noise level. Rather than fix it arbitrarily, we set
it by a method-of-moments empirical-Bayes pre-pass so that the prior matches the
dispersion of \emph{well-estimated} precisions:
\begin{enumerate}
  \item Run a short Gibbs pre-pass on the full training set with a non-informative,
  fixed noise prior ($a_0=1$, $b_0=1$, $b_0$ not learned) and collect the
  posterior draws $\{\alpha_c^{(s)}\}$.
  \item Restrict to data-rich chemicals ($N_c\ge N_{\min}$, with $N_{\min}=50$),
  whose posterior precisions are essentially data-driven, and form their
  posterior means $\hat{\alpha}_c=\tfrac{1}{S}\sum_s\alpha_c^{(s)}$.
  \item Fit a Gamma to $\{\hat{\alpha}_c\}$ by matching mean $m$ and variance
  $\nu$, and take the shape
  \begin{equation}
    a_0 \;=\; m^2/\nu.
  \end{equation}
  The corresponding rate is discarded, because the main run learns $b_0$ from its
  hyperprior (Eq.~\ref{eq:post-b0}).
\end{enumerate}
For the run reported in the paper this gives $a_0\approx1.75$; the main-run
posterior for the shared rate concentrates near $b_0\approx0.08$, so
$\Gam(a_0,b_0)$ is a weakly informative population prior on the per-chemical
precisions. The same EB $a_0$ is used for the cross-validation runs and for the
full-data model used to generate SSDs.

% =====================================================================
\subsection{Hyperparameter settings}
\label{sec:hyper}

Table~\ref{tab:hyper} lists the fixed constants and run settings. The mean- and
weight/factor-precision hyperpriors are the same trivial constants used by
\citet{rendle2011bayesian}; with $p$ in the thousands their influence on the
posterior is negligible. The noise-model hyperprior $\Gam(c_0,d_0)$ on $b_0$ and
the EB shape $a_0$ are specific to this work.

\begin{table}[htbp]
\centering
\caption{Hyperparameter and run settings.}
\label{tab:hyper}
\small
\begin{tabular}{lll}
\toprule
\textbf{Quantity} & \textbf{Symbol} & \textbf{Value} \\
\midrule
Factorization rank            & $k$                       & $32$ \\
Mean-hyperprior mean          & $\mu_0$                   & $0$ \\
Mean-hyperprior precision     & $\gamma_0$                & $1$ \\
Weight/factor precision prior & $\alpha_\lambda,\beta_\lambda$ & $1,\ 1$ \\
Noise-rate hyperprior         & $c_0,\ d_0$               & $1,\ 1$ \\
Noise-prior shape (empirical Bayes) & $a_0$               & $\approx 1.75$ (fixed after EB) \\
Noise-prior rate              & $b_0$                     & initialised at $1$, learned \\
Main run                      & $n_{\mathrm{iter}},\ n_{\mathrm{burn}}$ & $2000,\ 100$ \\
Retained samples              & $S$                       & $1900$ \\
EB pre-pass                   & $n_{\mathrm{iter}},\ n_{\mathrm{burn}}$ & $100,\ 50$ \\
EB data-rich threshold        & $N_{\min}$                & $50$ \\
Cross-validation              & folds                     & $5$ (grouped by triplet) \\
\bottomrule
\end{tabular}
\end{table}

% =====================================================================
\subsection{Posterior-predictive quantities}
\label{sec:predictive}

For a query $\xx^\ast$ belonging to chemical $c^\ast$, each retained sample gives
a deterministic FM prediction $\yhat^{(s)}(\xx^\ast)$ via Eq.~\eqref{eq:fm}. The
point prediction is the sample mean
\begin{equation}
  \yhat^\ast=\frac{1}{S}\sum_{s=1}^{S}\yhat^{(s)}(\xx^\ast).
\end{equation}
The predictive variance decomposes by the law of total variance,
$\Var(y^\ast\mid\mathcal{D})=\Var_\theta\!\big(\E[y^\ast\mid\theta]\big)
+\E_\theta\!\big[\Var(y^\ast\mid\theta)\big]$, whose two Monte-Carlo estimates are
the epistemic and aleatoric components:
\begin{align}
  \Var_{\mathrm{epi}}(\xx^\ast) &= \frac{1}{S}\sum_{s=1}^{S}\big(\yhat^{(s)}(\xx^\ast)-\yhat^\ast\big)^2,
  \label{eq:epi}\\
  \Var_{\mathrm{ale}}(c^\ast) &= \frac{1}{S}\sum_{s=1}^{S}\frac{1}{\alpha_{c^\ast}^{(s)}},
  \label{eq:ale}\\
  \Var_{\mathrm{tot}}(\xx^\ast,c^\ast) &= \Var_{\mathrm{epi}}(\xx^\ast)+\Var_{\mathrm{ale}}(c^\ast).
  \label{eq:tot}
\end{align}
The epistemic term is the spread of the mean across plausible fits and shrinks as
data constrain the parameters; the aleatoric term is the posterior-mean noise
variance of the chemical and does not vanish with more data. The aleatoric term
depends only on the chemical, so it is constant across species and durations for
a given chemical.

\paragraph{Calibration.} To assess whether $\Var_{\mathrm{tot}}$ is well
calibrated we form posterior-predictive draws that combine both components. For
each held-out observation and each sample $s$ we draw
\begin{equation}
  \tilde{y}^{(s)} \sim \Norm\!\big(\yhat^{(s)}(\xx^\ast),\ 1/\alpha_{c^\ast}^{(s)}\big),
\end{equation}
giving an $S$-point predictive sample. For a nominal central level $\rho$ we take
the empirical $[(1-\rho)/2,\ (1+\rho)/2]$ quantiles of these draws as the
credible interval, and report empirical coverage as the fraction of held-out
observations whose true value lies inside. Coverage is evaluated on the
out-of-fold observations (Section~\ref{sec:cv}), so predictions are scored by
models that did not train on them.

% =====================================================================
\subsection{SSD ensembles and hazardous concentrations}
\label{sec:ssd}

The Bayesian structure lets uncertainty propagate coherently into Species
Sensitivity Distributions. For a target chemical $c^\ast$ at a fixed duration,
and each retained sample $s$, we evaluate Eq.~\eqref{eq:fm} at $\theta^{(s)}$ for
all species $i=1,\dots,M$ and sort the resulting predictions ascending to obtain
one SSD curve $\big(\yhat^{(s)}_{(1)}\le\cdots\le\yhat^{(s)}_{(M)}\big)$ plotted
against the plotting positions $m/(M+1)$. Because every curve comes from a single
coherent parameter draw, the correlated epistemic structure across species is
preserved: the whole curve shifts and rescales as the shared parameters vary,
rather than species scattering independently. The ensemble of $S$ curves is
summarised pointwise by its median and central $95\%$ band.

A hazardous concentration $\mathrm{HC}x$ (the concentration protective of
$100-x\%$ of species) is read off each sampled curve at rank
$\lfloor x\,M/100\rfloor$, giving a posterior sample of $\mathrm{HC}x$ values;
its median and $2.5/97.5$ percentiles are the reported estimate and $95\%$
credible interval. Optionally, aleatoric noise can be injected per species by
adding a draw $\Norm\!\big(0,1/\alpha_{c^\ast}^{(s)}\big)$ to each prediction
before sorting; the SSD figures and HC intervals in the main text use the
epistemic ensemble (no added aleatoric noise) unless stated otherwise.

% =====================================================================
\subsection{Cross-validation and out-of-fold uncertainty}
\label{sec:cv}

Predictive accuracy and calibration are evaluated with grouped $5$-fold
cross-validation. Folds are formed by \emph{grouped} splitting on the
(chemical, species, duration) triplet identifier, so that all replicate
measurements of a held-out triplet are removed from training together and no
replicate of a scored triplet is ever seen during training. Each fold refits the
full model (including the hierarchical noise model, with the EB shape $a_0$ fixed
and $b_0$ learned) on its training partition and predicts the held-out partition;
concatenating the held-out predictions over folds yields the out-of-fold arrays
used for the RMSE, the uncertainty-versus-data analyses
(Eqs.~\ref{eq:epi}--\ref{eq:ale}), and the calibration curve. The full-data model
used to generate the SSDs and the all-chemical HC comparison is fit once on all
$n$ observations with the same settings.

Two practical points. First, the epistemic estimate of Eq.~\eqref{eq:epi} is only
meaningful for chemicals that appear in the training partition of a fold;
chemicals with fewer total observations than the number of folds can have folds
in which none of their rows are present, leaving their chemical-specific
parameters at the prior and deflating their apparent epistemic variance. Such
cold-start chemicals are excluded from the per-chemical
uncertainty-versus-data analyses in the main text. Second, the aleatoric
out-of-fold value for an observation is the sample mean of
$1/\alpha_{c}^{(s)}$ over that fold's retained samples, so it reflects the noise
precision learned from the fold's training data only.

% =====================================================================
\subsection{Computational complexity}
\label{sec:complexity}

Every conditional in Section~\ref{sec:sampler} is a single Gaussian or Gamma
draw, and the residual and factor caches make each parameter update a rank-one
correction touching only the non-zeros of the relevant feature column. One full
sweep therefore costs $O(k\,N_{\mathrm{nz}})$, linear in the factorization rank
and in the number of non-zero design-matrix entries, with an additional
$O(C)$ for the noise block (Eqs.~\ref{eq:post-alpha}--\ref{eq:post-b0}). This
matches the complexity of the unblocked BFM sampler of
\citet{rendle2011bayesian}; the per-chemical noise model adds only the vectorised
$\alpha_c$ draw and one scalar $b_0$ draw per sweep, so the hierarchical
extension does not change the asymptotic cost.

.
% Relies on the preamble of main.tex (amsmath, booktabs, etc. are already
% loaded there; amssymb is added there for \Real below). Uses the shared
% references.bib already cited by the rest of the paper.
% =====================================================================

% ---- convenience macros ----
\newcommand{\yhat}{\hat{y}}
\newcommand{\Norm}{\mathcal{N}}
\newcommand{\Gam}{\mathrm{Gamma}}
\newcommand{\Real}{\mathbb{R}}
\newcommand{\E}{\mathbb{E}}
\newcommand{\Var}{\mathrm{Var}}
\newcommand{\vv}{\mathbf{v}}
\newcommand{\xx}{\mathbf{x}}

\section{Supporting Information}
\label{sec:si}

This Supporting Information gives the full mathematical specification of the
Hierarchical Bayesian Factorization Machine (BFM) summarised in the main text.
It is written for a statistical or machine-learning audience and documents the
model exactly as implemented: the predictor, the complete hierarchical
generative model, the Gibbs sampler and every conditional posterior, the
empirical-Bayes initialisation, the hyperparameter settings, and the
posterior-predictive quantities used for calibration, Species Sensitivity
Distributions (SSDs), and hazardous concentrations. The predictive core is the
factorization machine (FM) of \citet{rendle2010fm} cast in the Bayesian form of
\citet{rendle2011bayesian}; Section~\ref{sec:diff} states precisely how our
model differs from theirs. Symbols follow the main text and are collected in
Table~\ref{tab:notation}.

% =====================================================================
\subsection{Notation}
\label{sec:notation}

\begin{table}[htbp]
\centering
\caption{Notation used throughout the Supporting Information.}
\label{tab:notation}
\small
\begin{tabular}{ll}
\toprule
\textbf{Symbol} & \textbf{Meaning} \\
\midrule
$n$ & number of observations (rows) \\
$p$ & number of features (design-matrix columns) \\
$k$ & number of latent factors per feature (factorization rank) \\
$C$ & number of chemicals (noise groups) \\
$\xx_i \in \Real^{p}$ & feature vector of observation $i$ \\
$x_{ij}$ & $j$-th entry of $\xx_i$ \\
$y_i$ & centred log-transformed LC50 of observation $i$ \\
$c_i \in \{1,\dots,C\}$ & chemical (group) index of observation $i$ \\
$w_0$ & global bias \\
$w_j$ & linear weight of feature $j$ \\
$v_{jf}$ & $f$-th latent factor of feature $j$; $\vv_j=(v_{j1},\dots,v_{jk})$ \\
$\alpha_c$ & noise precision of chemical $c$ (aleatoric term) \\
$\mu_w,\lambda_w$ & mean and precision of the shared prior on the weights $w_j$ \\
$\mu_v,\lambda_v$ & mean and precision of the shared prior on the factors $v_{jf}$ \\
$a_0,b_0$ & shape and rate of the shared prior on the $\alpha_c$ \\
$\mu_0,\gamma_0$ & mean and precision of the Gaussian mean-hyperpriors \\
$\alpha_\lambda,\beta_\lambda$ & shape and rate of the Gamma prior on $\lambda_w,\lambda_v$ \\
$c_0,d_0$ & shape and rate of the Gamma hyperprior on $b_0$ \\
$\theta$ & all model parameters $\{w_0,\{w_j\},\{v_{jf}\}\}$ \\
$S$ & number of retained (post-burn-in) posterior samples \\
$\yhat^{(s)}$ & prediction from posterior sample $s$ \\
\bottomrule
\end{tabular}
\end{table}

We write $\Norm(\mu,\sigma^2)$ for a Gaussian with mean $\mu$ and variance
$\sigma^2$, and $\Gam(a,b)$ for a Gamma distribution with shape $a$ and rate $b$
(so that its mean is $a/b$). All precisions are inverse variances.

% =====================================================================
\subsection{Data representation and feature encoding}
\label{sec:features}

Each observation is a single laboratory LC50 measurement for a
(chemical, species, duration) combination. The target $y_i$ is the
log-transformed concentration in mg/L, centred by subtracting the global mean
$\bar{y}$ of the training set; the mean is added back to recover log-mg/L units
when reporting predictions. Repeated measurements of the same triplet are kept
as separate rows rather than averaged, so that the per-chemical noise model
(Section~\ref{sec:model}) sees the within-triplet spread directly.

The feature vector $\xx_i$ concatenates five one-hot categorical blocks and two
standardised numerical features:
\begin{itemize}
  \item \textbf{chemical} identity (one-hot over $C$ chemicals),
  \item \textbf{species} identity (one-hot),
  \item \textbf{test duration} (one-hot over $\{24,48,72,96\}$\,h),
  \item \textbf{taxonomic family} and \textbf{taxonomic class} of the species (two one-hot blocks),
  \item \textbf{numerical}: log-transformed molecular weight and the RDKit
        octanol--water partition coefficient (\texttt{clogp}), median-imputed
        and standardised to zero mean and unit variance.
\end{itemize}
Unseen categories at prediction time map to an all-zero block (they contribute
nothing to the sums below), so the model can score any (chemical, species,
duration) triplet whose chemical and species were each observed at least once
during training, but not entirely novel chemicals or species. For the dataset
used here $n=70{,}670$, $C=3{,}295$ chemicals, $1{,}267$ species, and
$p=4{,}879$ (namely $3{,}295+1{,}267+4+283+28+2$).

% =====================================================================
\subsection{The factorization-machine predictor}
\label{sec:fm}

The predictive mean for observation $i$ is the second-order factorization
machine
\begin{equation}
  \yhat_i(\theta) \;=\; w_0 \;+\; \sum_{j=1}^{p} w_j\,x_{ij}
  \;+\; \sum_{j=1}^{p}\sum_{j'>j}^{p} \langle \vv_j,\vv_{j'}\rangle\,x_{ij}x_{ij'},
  \label{eq:fm}
\end{equation}
with $\langle \vv_j,\vv_{j'}\rangle=\sum_{f=1}^{k} v_{jf}v_{j'f}$. The pairwise
interactions are not stored explicitly; they are represented in the low-rank
form of Eq.~\eqref{eq:fm}, which lets the model infer the interaction between a
chemical and a species that were never tested together. Using the standard FM
identity, the interaction sum is evaluated in $O(pk)$ time (in practice $O(k
N_{\mathrm{nz}})$ for a sparse $\xx_i$ with $N_{\mathrm{nz}}$ non-zeros) as
\begin{equation}
  \sum_{j<j'} \langle \vv_j,\vv_{j'}\rangle\,x_{ij}x_{ij'}
  \;=\; \tfrac{1}{2}\sum_{f=1}^{k}\!\left[
      \Big(\sum_{j=1}^{p} v_{jf}x_{ij}\Big)^{\!2}
      - \sum_{j=1}^{p} v_{jf}^{2}x_{ij}^{2}
  \right].
  \label{eq:fm-identity}
\end{equation}
The diagonal subtraction in Eq.~\eqref{eq:fm-identity} removes each feature's
self-interaction. A direct consequence, used repeatedly by the sampler, is that
$\yhat_i$ is an \emph{affine} function of each individual parameter when the
others are held fixed (Section~\ref{sec:sampler}).

% =====================================================================
\subsection{Full hierarchical generative model}
\label{sec:model}

The model is a Bayesian FM in which the single global noise precision of the
standard BFM is replaced by a per-chemical precision $\alpha_c$ under a shared
Gamma prior. The complete generative model is:

\paragraph{Likelihood.} Each observation is its FM mean plus Gaussian noise
whose precision is that of the observation's chemical,
\begin{equation}
  y_i \;\mid\; \xx_i,\theta,\alpha_{c_i}
  \;\sim\; \Norm\!\big(\yhat_i(\theta),\ \alpha_{c_i}^{-1}\big),
  \qquad i=1,\dots,n.
  \label{eq:like}
\end{equation}

\paragraph{Priors on the model parameters.} The global bias has a fixed Gaussian
prior; the weights and factors are drawn from shared populations whose mean and
precision are themselves learned:
\begin{align}
  w_0     &\sim \Norm(\mu_0,\ \gamma_0^{-1}), \\
  w_j     &\sim \Norm(\mu_w,\ \lambda_w^{-1}), & j&=1,\dots,p, \label{eq:prior-w}\\
  v_{jf}  &\sim \Norm(\mu_v,\ \lambda_v^{-1}), & j&=1,\dots,p,\ f=1,\dots,k. \label{eq:prior-v}
\end{align}

\paragraph{Hyperpriors on the populations.} There is one population for all
linear weights and one for all latent factors (as in the base BFM of
\citet{rendle2011bayesian}); each has a Gaussian hyperprior on its mean and a
Gamma hyperprior on its precision:
\begin{align}
  \mu_w &\sim \Norm(\mu_0,\ \gamma_0^{-1}), & \lambda_w &\sim \Gam(\alpha_\lambda,\ \beta_\lambda), \\
  \mu_v &\sim \Norm(\mu_0,\ \gamma_0^{-1}), & \lambda_v &\sim \Gam(\alpha_\lambda,\ \beta_\lambda).
\end{align}

\paragraph{Hierarchical noise model.} This is the methodological contribution.
Each chemical has its own noise precision, drawn from a Gamma prior whose rate is
shared across all chemicals and is itself given a Gamma hyperprior:
\begin{align}
  \alpha_c &\sim \Gam(a_0,\ b_0), & c&=1,\dots,C, \label{eq:prior-alpha}\\
  b_0      &\sim \Gam(c_0,\ d_0). \label{eq:prior-b0}
\end{align}
The shape $a_0$ is held fixed (set by empirical Bayes, Section~\ref{sec:eb});
the rate $b_0$ is resampled every sweep. Because $b_0$ is common to all
chemicals, it ties the per-chemical precisions together: a chemical with many
observations has its $\alpha_c$ set mainly by the scatter of its own residuals,
whereas a sparsely measured chemical is pooled towards the population-typical
noise level implied by $\Gam(a_0,b_0)$. This partial pooling is what keeps the
aleatoric estimate stable for data-poor chemicals.

The joint density factorises as
\begin{equation}
\begin{aligned}
  p\big(\theta,\{\alpha_c\},b_0,\mu_w,\lambda_w,\mu_v,\lambda_v \mid \mathbf{y}\big)
  \;\propto\;
  & \Big[\textstyle\prod_{i=1}^{n} \Norm\!\big(y_i \mid \yhat_i(\theta),\alpha_{c_i}^{-1}\big)\Big]
    \Norm(w_0\mid\mu_0,\gamma_0^{-1}) \\
  &\times \Big[\textstyle\prod_{j=1}^{p}\Norm(w_j\mid\mu_w,\lambda_w^{-1})\Big]
          \Big[\textstyle\prod_{j,f}\Norm(v_{jf}\mid\mu_v,\lambda_v^{-1})\Big] \\
  &\times \Norm(\mu_w\mid\mu_0,\gamma_0^{-1})\,\Gam(\lambda_w\mid\alpha_\lambda,\beta_\lambda)\,
          \Norm(\mu_v\mid\mu_0,\gamma_0^{-1})\,\Gam(\lambda_v\mid\alpha_\lambda,\beta_\lambda) \\
  &\times \Big[\textstyle\prod_{c=1}^{C}\Gam(\alpha_c\mid a_0,b_0)\Big]\,\Gam(b_0\mid c_0,d_0).
\end{aligned}
\label{eq:joint}
\end{equation}
This is the model drawn in plate notation in the main text: the fixed constants
$(\mu_0,\gamma_0,\alpha_\lambda,\beta_\lambda,a_0,c_0,d_0)$ set the hyperpriors;
$w_j$ and $v_{jf}$ share $(\mu_w,\lambda_w)$ and $(\mu_v,\lambda_v)$
respectively; the $\alpha_c$ share the rate $b_0$; and each $y_i$ is drawn from a
Gaussian centred at $\yhat_i$ with precision $\alpha_{c_i}$.

% =====================================================================
\subsection{Relationship to the standard Bayesian FM}
\label{sec:diff}

Our model departs from the BFM of \citet{rendle2011bayesian} in two respects.
\begin{enumerate}
  \item \textbf{Per-chemical noise (the substantive change).} The standard BFM
  assumes a single global precision $\alpha$ shared by all observations, with a
  Gamma prior and a Gamma-conjugate posterior. We replace it by $C$
  chemical-specific precisions $\alpha_c$ (Eq.~\ref{eq:prior-alpha}) under a
  shared, learned rate $b_0$ (Eq.~\ref{eq:prior-b0}). Equivalently, the global
  $\alpha$ in every conditional below becomes the observation-specific
  $\alpha_{c_i}$, and one extra conjugate draw for $b_0$ is added per sweep. This
  turns the noise term into a per-chemical, partially pooled quantity, which is
  the aleatoric component carried into the downstream analysis.

  \item \textbf{Independent Gaussian priors on the population means (a minor
  change).} The standard BFM couples each population mean to its precision
  through a Normal--Gamma prior, $\mu \mid \lambda \sim
  \Norm(\mu_0,(\gamma_0\lambda)^{-1})$, so the posterior precision of $\mu$
  scales with $\lambda$. We instead place an independent
  $\Norm(\mu_0,\gamma_0^{-1})$ prior on each mean (Eqs.~\ref{eq:prior-w}--\ref{eq:prior-v}),
  decoupling $\mu$ from $\lambda$; the mean does not then enter the update for
  $\lambda$, and the $\mu$-posterior precision is $\gamma_0+p\lambda$ rather than
  $(\gamma_0+p)\lambda$ (Section~\ref{sec:sampler}). With $p$ in the thousands
  this has negligible effect on inference, as \citet{rendle2011bayesian} also
  note for the hyperprior constants, but we state it for exactness.
\end{enumerate}
As in the reference, we use single-parameter (unblocked) Gibbs sampling, which
avoids the matrix inversions of blocked schemes and keeps each sweep linear in
$k$.

% =====================================================================
\subsection{Gibbs sampler and conditional posteriors}
\label{sec:sampler}

Posterior inference is by Gibbs sampling: each quantity is drawn in turn from its
full conditional given the current values of the others. All conditionals are
available in closed form because (i) the FM mean is affine in each single
parameter, giving Gaussian conditionals with Gaussian priors, and (ii) the
Gamma priors on the precisions are conjugate to the Gaussian likelihood.

Throughout we cache the running residual
\begin{equation}
  e_i \;=\; y_i - \yhat_i(\theta),
  \label{eq:resid}
\end{equation}
and the per-factor cache $q_{if}=\sum_{j=1}^{p} v_{jf}x_{ij}$, both updated by a
rank-one correction after each parameter is resampled. This is what makes a full
sweep $O(k N_{\mathrm{nz}})$, where $N_{\mathrm{nz}}$ is the number of non-zeros
in the design matrix. For a parameter with current value $\theta_\ell$ and
``slope'' $h_{i\ell}=\partial \yhat_i/\partial\theta_\ell$ (constant in
$\theta_\ell$ by multilinearity), the partial residual that adds back its current
contribution is $r_{i}=e_i+h_{i\ell}\theta_\ell$.

\subsubsection{Global bias $w_0$}
Here $h_{i0}=1$ for all $i$, so with prior $\Norm(\mu_0,\gamma_0^{-1})$,
\begin{equation}
  \sigma_{w_0}^2=\Big(\gamma_0+\sum_{i=1}^{n}\alpha_{c_i}\Big)^{-1},
  \qquad
  \mu_{w_0}=\sigma_{w_0}^2\Big(\gamma_0\mu_0+\sum_{i=1}^{n}\alpha_{c_i}\,(e_i+w_0)\Big),
  \qquad
  w_0 \sim \Norm(\mu_{w_0},\sigma_{w_0}^2).
\end{equation}
Then set $e_i \mathrel{+}= (w_0^{\mathrm{old}}-w_0^{\mathrm{new}})$ for all $i$.

\subsubsection{Weight population $(\mu_w,\lambda_w)$}
Treating the $p$ weights as i.i.d.\ $\Norm(\mu_w,\lambda_w^{-1})$ draws,
\begin{align}
  \mu_w &\sim \Norm\!\big(\mu_{\mu_w},\,\sigma_{\mu_w}^2\big),
  & \sigma_{\mu_w}^2 &=\big(\gamma_0+p\,\lambda_w\big)^{-1},
  & \mu_{\mu_w} &=\sigma_{\mu_w}^2\big(\gamma_0\mu_0+\lambda_w\textstyle\sum_{j=1}^{p}w_j\big), \\
  \lambda_w &\sim \Gam\!\Big(\alpha_\lambda+\tfrac{p}{2},\ \
     \beta_\lambda+\tfrac{1}{2}\textstyle\sum_{j=1}^{p}(w_j-\mu_w)^2\Big).
\end{align}

\subsubsection{Linear weights $w_j$}
For feature $j$, $h_{ij}=x_{ij}$, and only observations with $x_{ij}\neq 0$
contribute. With prior $\Norm(\mu_w,\lambda_w^{-1})$ and partial residual
$r_i=e_i+x_{ij}w_j$,
\begin{equation}
  \sigma_{w_j}^2=\Big(\lambda_w+\sum_{i:\,x_{ij}\neq0}\!\alpha_{c_i}\,x_{ij}^2\Big)^{-1},
  \quad
  \mu_{w_j}=\sigma_{w_j}^2\Big(\lambda_w\mu_w+\sum_{i:\,x_{ij}\neq0}\!\alpha_{c_i}\,x_{ij}\,r_i\Big),
  \quad
  w_j\sim\Norm(\mu_{w_j},\sigma_{w_j}^2).
\end{equation}
Then update $e_i \mathrel{+}= x_{ij}\,(w_j^{\mathrm{old}}-w_j^{\mathrm{new}})$.

\subsubsection{Factor population $(\mu_v,\lambda_v)$}
Treating all $pk$ factor entries as i.i.d.\ $\Norm(\mu_v,\lambda_v^{-1})$,
\begin{align}
  \mu_v &\sim \Norm\!\big(\mu_{\mu_v},\,\sigma_{\mu_v}^2\big),
  & \sigma_{\mu_v}^2 &=\big(\gamma_0+p\,k\,\lambda_v\big)^{-1},
  & \mu_{\mu_v} &=\sigma_{\mu_v}^2\big(\gamma_0\mu_0+\lambda_v\textstyle\sum_{j,f}v_{jf}\big), \\
  \lambda_v &\sim \Gam\!\Big(\alpha_\lambda+\tfrac{p k}{2},\ \
     \beta_\lambda+\tfrac{1}{2}\textstyle\sum_{j,f}(v_{jf}-\mu_v)^2\Big).
\end{align}

\subsubsection{Latent factors $v_{jf}$}
By the multilinearity of Eq.~\eqref{eq:fm-identity}, $\yhat_i$ is affine in
$v_{jf}$ with slope
\begin{equation}
  h_{ijf}=x_{ij}\big(q_{if}-x_{ij}v_{jf}\big),
  \qquad q_{if}=\sum_{j'=1}^{p} v_{j'f}x_{ij'},
\end{equation}
which does not depend on $v_{jf}$ (the $x_{ij}v_{jf}$ subtraction removes it).
With prior $\Norm(\mu_v,\lambda_v^{-1})$ and partial residual
$r_i=e_i+h_{ijf}v_{jf}$,
\begin{equation}
  \sigma_{v_{jf}}^2=\Big(\lambda_v+\sum_{i:\,x_{ij}\neq0}\!\alpha_{c_i}\,h_{ijf}^2\Big)^{-1},
  \quad
  \mu_{v_{jf}}=\sigma_{v_{jf}}^2\Big(\lambda_v\mu_v+\sum_{i:\,x_{ij}\neq0}\!\alpha_{c_i}\,h_{ijf}\,r_i\Big),
  \quad
  v_{jf}\sim\Norm(\mu_{v_{jf}},\sigma_{v_{jf}}^2).
\end{equation}
Then update both caches: $e_i \mathrel{+}= h_{ijf}\,(v_{jf}^{\mathrm{old}}-v_{jf}^{\mathrm{new}})$
and $q_{if}\mathrel{+}= x_{ij}\,(v_{jf}^{\mathrm{new}}-v_{jf}^{\mathrm{old}})$.

\subsubsection{Per-chemical noise precisions $\alpha_c$}
Let $N_c=|\{i:c_i=c\}|$ and $\mathrm{SSE}_c=\sum_{i:\,c_i=c} e_i^2$ be the
observation count and residual sum of squares of chemical $c$. Conjugacy of the
Gamma prior (Eq.~\ref{eq:prior-alpha}) to the Gaussian likelihood
(Eq.~\ref{eq:like}) gives
\begin{equation}
  \alpha_c \;\sim\; \Gam\!\Big(a_0+\tfrac{N_c}{2},\ \ b_0+\tfrac{1}{2}\,\mathrm{SSE}_c\Big),
  \qquad c=1,\dots,C.
  \label{eq:post-alpha}
\end{equation}
All $C$ precisions are drawn in one vectorised step using per-group sums of
$e_i^2$.

\subsubsection{Shared noise rate $b_0$}
The rate $b_0$ is conjugate to the $C$ Gamma-distributed precisions with common
shape $a_0$: combining $\prod_c \Gam(\alpha_c\mid a_0,b_0)\propto
b_0^{C a_0}e^{-b_0\sum_c\alpha_c}$ with the prior $\Gam(c_0,d_0)$ yields
\begin{equation}
  b_0 \;\sim\; \Gam\!\Big(c_0+C\,a_0,\ \ d_0+\textstyle\sum_{c=1}^{C}\alpha_c\Big).
  \label{eq:post-b0}
\end{equation}

\subsubsection{One sweep and sample collection}
A single Gibbs sweep applies, in order: $w_0$; $(\mu_w,\lambda_w)$ then every
$w_j$; $(\mu_v,\lambda_v)$ then every $v_{jf}$; the $\alpha_c$
(Eq.~\ref{eq:post-alpha}); and $b_0$ (Eq.~\ref{eq:post-b0}). The chain is run
for $n_{\mathrm{iter}}$ sweeps from an arbitrary initialisation
($w_0=0$, $w=\mathbf{0}$, $v_{jf}\sim\Norm(0,0.1^2)$, $\alpha_c=1$); the first
$n_{\mathrm{burn}}$ sweeps are discarded as burn-in and the remaining
$S=n_{\mathrm{iter}}-n_{\mathrm{burn}}$ sweeps are retained. Each retained sample
$s$ stores $\{w_0^{(s)},w^{(s)},v^{(s)},\{\alpha_c^{(s)}\},b_0^{(s)}\}$ and
corresponds to one complete, internally consistent parameterisation of the
model.

% =====================================================================
\subsection{Empirical-Bayes initialisation of the prior shape $a_0$}
\label{sec:eb}

The shared shape $a_0$ controls how strongly sparsely tested chemicals are
pooled towards the population noise level. Rather than fix it arbitrarily, we set
it by a method-of-moments empirical-Bayes pre-pass so that the prior matches the
dispersion of \emph{well-estimated} precisions:
\begin{enumerate}
  \item Run a short Gibbs pre-pass on the full training set with a non-informative,
  fixed noise prior ($a_0=1$, $b_0=1$, $b_0$ not learned) and collect the
  posterior draws $\{\alpha_c^{(s)}\}$.
  \item Restrict to data-rich chemicals ($N_c\ge N_{\min}$, with $N_{\min}=50$),
  whose posterior precisions are essentially data-driven, and form their
  posterior means $\hat{\alpha}_c=\tfrac{1}{S}\sum_s\alpha_c^{(s)}$.
  \item Fit a Gamma to $\{\hat{\alpha}_c\}$ by matching mean $m$ and variance
  $\nu$, and take the shape
  \begin{equation}
    a_0 \;=\; m^2/\nu.
  \end{equation}
  The corresponding rate is discarded, because the main run learns $b_0$ from its
  hyperprior (Eq.~\ref{eq:post-b0}).
\end{enumerate}
For the run reported in the paper this gives $a_0\approx1.75$; the main-run
posterior for the shared rate concentrates near $b_0\approx0.08$, so
$\Gam(a_0,b_0)$ is a weakly informative population prior on the per-chemical
precisions. The same EB $a_0$ is used for the cross-validation runs and for the
full-data model used to generate SSDs.

% =====================================================================
\subsection{Hyperparameter settings}
\label{sec:hyper}

Table~\ref{tab:hyper} lists the fixed constants and run settings. The mean- and
weight/factor-precision hyperpriors are the same trivial constants used by
\citet{rendle2011bayesian}; with $p$ in the thousands their influence on the
posterior is negligible. The noise-model hyperprior $\Gam(c_0,d_0)$ on $b_0$ and
the EB shape $a_0$ are specific to this work.

\begin{table}[htbp]
\centering
\caption{Hyperparameter and run settings.}
\label{tab:hyper}
\small
\begin{tabular}{lll}
\toprule
\textbf{Quantity} & \textbf{Symbol} & \textbf{Value} \\
\midrule
Factorization rank            & $k$                       & $32$ \\
Mean-hyperprior mean          & $\mu_0$                   & $0$ \\
Mean-hyperprior precision     & $\gamma_0$                & $1$ \\
Weight/factor precision prior & $\alpha_\lambda,\beta_\lambda$ & $1,\ 1$ \\
Noise-rate hyperprior         & $c_0,\ d_0$               & $1,\ 1$ \\
Noise-prior shape (empirical Bayes) & $a_0$               & $\approx 1.75$ (fixed after EB) \\
Noise-prior rate              & $b_0$                     & initialised at $1$, learned \\
Main run                      & $n_{\mathrm{iter}},\ n_{\mathrm{burn}}$ & $2000,\ 100$ \\
Retained samples              & $S$                       & $1900$ \\
EB pre-pass                   & $n_{\mathrm{iter}},\ n_{\mathrm{burn}}$ & $100,\ 50$ \\
EB data-rich threshold        & $N_{\min}$                & $50$ \\
Cross-validation              & folds                     & $5$ (grouped by triplet) \\
\bottomrule
\end{tabular}
\end{table}

% =====================================================================
\subsection{Posterior-predictive quantities}
\label{sec:predictive}

For a query $\xx^\ast$ belonging to chemical $c^\ast$, each retained sample gives
a deterministic FM prediction $\yhat^{(s)}(\xx^\ast)$ via Eq.~\eqref{eq:fm}. The
point prediction is the sample mean
\begin{equation}
  \yhat^\ast=\frac{1}{S}\sum_{s=1}^{S}\yhat^{(s)}(\xx^\ast).
\end{equation}
The predictive variance decomposes by the law of total variance,
$\Var(y^\ast\mid\mathcal{D})=\Var_\theta\!\big(\E[y^\ast\mid\theta]\big)
+\E_\theta\!\big[\Var(y^\ast\mid\theta)\big]$, whose two Monte-Carlo estimates are
the epistemic and aleatoric components:
\begin{align}
  \Var_{\mathrm{epi}}(\xx^\ast) &= \frac{1}{S}\sum_{s=1}^{S}\big(\yhat^{(s)}(\xx^\ast)-\yhat^\ast\big)^2,
  \label{eq:epi}\\
  \Var_{\mathrm{ale}}(c^\ast) &= \frac{1}{S}\sum_{s=1}^{S}\frac{1}{\alpha_{c^\ast}^{(s)}},
  \label{eq:ale}\\
  \Var_{\mathrm{tot}}(\xx^\ast,c^\ast) &= \Var_{\mathrm{epi}}(\xx^\ast)+\Var_{\mathrm{ale}}(c^\ast).
  \label{eq:tot}
\end{align}
The epistemic term is the spread of the mean across plausible fits and shrinks as
data constrain the parameters; the aleatoric term is the posterior-mean noise
variance of the chemical and does not vanish with more data. The aleatoric term
depends only on the chemical, so it is constant across species and durations for
a given chemical.

\paragraph{Calibration.} To assess whether $\Var_{\mathrm{tot}}$ is well
calibrated we form posterior-predictive draws that combine both components. For
each held-out observation and each sample $s$ we draw
\begin{equation}
  \tilde{y}^{(s)} \sim \Norm\!\big(\yhat^{(s)}(\xx^\ast),\ 1/\alpha_{c^\ast}^{(s)}\big),
\end{equation}
giving an $S$-point predictive sample. For a nominal central level $\rho$ we take
the empirical $[(1-\rho)/2,\ (1+\rho)/2]$ quantiles of these draws as the
credible interval, and report empirical coverage as the fraction of held-out
observations whose true value lies inside. Coverage is evaluated on the
out-of-fold observations (Section~\ref{sec:cv}), so predictions are scored by
models that did not train on them.

% =====================================================================
\subsection{SSD ensembles and hazardous concentrations}
\label{sec:ssd}

The Bayesian structure lets uncertainty propagate coherently into Species
Sensitivity Distributions. For a target chemical $c^\ast$ at a fixed duration,
and each retained sample $s$, we evaluate Eq.~\eqref{eq:fm} at $\theta^{(s)}$ for
all species $i=1,\dots,M$ and sort the resulting predictions ascending to obtain
one SSD curve $\big(\yhat^{(s)}_{(1)}\le\cdots\le\yhat^{(s)}_{(M)}\big)$ plotted
against the plotting positions $m/(M+1)$. Because every curve comes from a single
coherent parameter draw, the correlated epistemic structure across species is
preserved: the whole curve shifts and rescales as the shared parameters vary,
rather than species scattering independently. The ensemble of $S$ curves is
summarised pointwise by its median and central $95\%$ band.

A hazardous concentration $\mathrm{HC}x$ (the concentration protective of
$100-x\%$ of species) is read off each sampled curve at rank
$\lfloor x\,M/100\rfloor$, giving a posterior sample of $\mathrm{HC}x$ values;
its median and $2.5/97.5$ percentiles are the reported estimate and $95\%$
credible interval. Optionally, aleatoric noise can be injected per species by
adding a draw $\Norm\!\big(0,1/\alpha_{c^\ast}^{(s)}\big)$ to each prediction
before sorting; the SSD figures and HC intervals in the main text use the
epistemic ensemble (no added aleatoric noise) unless stated otherwise.

% =====================================================================
\subsection{Cross-validation and out-of-fold uncertainty}
\label{sec:cv}

Predictive accuracy and calibration are evaluated with grouped $5$-fold
cross-validation. Folds are formed by \emph{grouped} splitting on the
(chemical, species, duration) triplet identifier, so that all replicate
measurements of a held-out triplet are removed from training together and no
replicate of a scored triplet is ever seen during training. Each fold refits the
full model (including the hierarchical noise model, with the EB shape $a_0$ fixed
and $b_0$ learned) on its training partition and predicts the held-out partition;
concatenating the held-out predictions over folds yields the out-of-fold arrays
used for the RMSE, the uncertainty-versus-data analyses
(Eqs.~\ref{eq:epi}--\ref{eq:ale}), and the calibration curve. The full-data model
used to generate the SSDs and the all-chemical HC comparison is fit once on all
$n$ observations with the same settings.

Two practical points. First, the epistemic estimate of Eq.~\eqref{eq:epi} is only
meaningful for chemicals that appear in the training partition of a fold;
chemicals with fewer total observations than the number of folds can have folds
in which none of their rows are present, leaving their chemical-specific
parameters at the prior and deflating their apparent epistemic variance. Such
cold-start chemicals are excluded from the per-chemical
uncertainty-versus-data analyses in the main text. Second, the aleatoric
out-of-fold value for an observation is the sample mean of
$1/\alpha_{c}^{(s)}$ over that fold's retained samples, so it reflects the noise
precision learned from the fold's training data only.

% =====================================================================
\subsection{Computational complexity}
\label{sec:complexity}

Every conditional in Section~\ref{sec:sampler} is a single Gaussian or Gamma
draw, and the residual and factor caches make each parameter update a rank-one
correction touching only the non-zeros of the relevant feature column. One full
sweep therefore costs $O(k\,N_{\mathrm{nz}})$, linear in the factorization rank
and in the number of non-zero design-matrix entries, with an additional
$O(C)$ for the noise block (Eqs.~\ref{eq:post-alpha}--\ref{eq:post-b0}). This
matches the complexity of the unblocked BFM sampler of
\citet{rendle2011bayesian}; the per-chemical noise model adds only the vectorised
$\alpha_c$ draw and one scalar $b_0$ draw per sweep, so the hierarchical
extension does not change the asymptotic cost.

.
% Relies on the preamble of main.tex (amsmath, booktabs, etc. are already
% loaded there; amssymb is added there for \Real below). Uses the shared
% references.bib already cited by the rest of the paper.
% =====================================================================

% ---- convenience macros ----
\newcommand{\yhat}{\hat{y}}
\newcommand{\Norm}{\mathcal{N}}
\newcommand{\Gam}{\mathrm{Gamma}}
\newcommand{\Real}{\mathbb{R}}
\newcommand{\E}{\mathbb{E}}
\newcommand{\Var}{\mathrm{Var}}
\newcommand{\vv}{\mathbf{v}}
\newcommand{\xx}{\mathbf{x}}

\section{Supporting Information}
\label{sec:si}

This Supporting Information gives the full mathematical specification of the
Hierarchical Bayesian Factorization Machine (BFM) summarised in the main text.
It is written for a statistical or machine-learning audience and documents the
model exactly as implemented: the predictor, the complete hierarchical
generative model, the Gibbs sampler and every conditional posterior, the
empirical-Bayes initialisation, the hyperparameter settings, and the
posterior-predictive quantities used for calibration, Species Sensitivity
Distributions (SSDs), and hazardous concentrations. The predictive core is the
factorization machine (FM) of \citet{rendle2010fm} cast in the Bayesian form of
\citet{rendle2011bayesian}; Section~\ref{sec:diff} states precisely how our
model differs from theirs. Symbols follow the main text and are collected in
Table~\ref{tab:notation}.

% =====================================================================
\subsection{Notation}
\label{sec:notation}

\begin{table}[htbp]
\centering
\caption{Notation used throughout the Supporting Information.}
\label{tab:notation}
\small
\begin{tabular}{ll}
\toprule
\textbf{Symbol} & \textbf{Meaning} \\
\midrule
$n$ & number of observations (rows) \\
$p$ & number of features (design-matrix columns) \\
$k$ & number of latent factors per feature (factorization rank) \\
$C$ & number of chemicals (noise groups) \\
$\xx_i \in \Real^{p}$ & feature vector of observation $i$ \\
$x_{ij}$ & $j$-th entry of $\xx_i$ \\
$y_i$ & centred log-transformed LC50 of observation $i$ \\
$c_i \in \{1,\dots,C\}$ & chemical (group) index of observation $i$ \\
$w_0$ & global bias \\
$w_j$ & linear weight of feature $j$ \\
$v_{jf}$ & $f$-th latent factor of feature $j$; $\vv_j=(v_{j1},\dots,v_{jk})$ \\
$\alpha_c$ & noise precision of chemical $c$ (aleatoric term) \\
$\mu_w,\lambda_w$ & mean and precision of the shared prior on the weights $w_j$ \\
$\mu_v,\lambda_v$ & mean and precision of the shared prior on the factors $v_{jf}$ \\
$a_0,b_0$ & shape and rate of the shared prior on the $\alpha_c$ \\
$\mu_0,\gamma_0$ & mean and precision of the Gaussian mean-hyperpriors \\
$\alpha_\lambda,\beta_\lambda$ & shape and rate of the Gamma prior on $\lambda_w,\lambda_v$ \\
$c_0,d_0$ & shape and rate of the Gamma hyperprior on $b_0$ \\
$\theta$ & all model parameters $\{w_0,\{w_j\},\{v_{jf}\}\}$ \\
$S$ & number of retained (post-burn-in) posterior samples \\
$\yhat^{(s)}$ & prediction from posterior sample $s$ \\
\bottomrule
\end{tabular}
\end{table}

We write $\Norm(\mu,\sigma^2)$ for a Gaussian with mean $\mu$ and variance
$\sigma^2$, and $\Gam(a,b)$ for a Gamma distribution with shape $a$ and rate $b$
(so that its mean is $a/b$). All precisions are inverse variances.

% =====================================================================
\subsection{Data representation and feature encoding}
\label{sec:features}

Each observation is a single laboratory LC50 measurement for a
(chemical, species, duration) combination. The target $y_i$ is the
log-transformed concentration in mg/L, centred by subtracting the global mean
$\bar{y}$ of the training set; the mean is added back to recover log-mg/L units
when reporting predictions. Repeated measurements of the same triplet are kept
as separate rows rather than averaged, so that the per-chemical noise model
(Section~\ref{sec:model}) sees the within-triplet spread directly.

The feature vector $\xx_i$ concatenates five one-hot categorical blocks and two
standardised numerical features:
\begin{itemize}
  \item \textbf{chemical} identity (one-hot over $C$ chemicals),
  \item \textbf{species} identity (one-hot),
  \item \textbf{test duration} (one-hot over $\{24,48,72,96\}$\,h),
  \item \textbf{taxonomic family} and \textbf{taxonomic class} of the species (two one-hot blocks),
  \item \textbf{numerical}: log-transformed molecular weight and the RDKit
        octanol--water partition coefficient (\texttt{clogp}), median-imputed
        and standardised to zero mean and unit variance.
\end{itemize}
Unseen categories at prediction time map to an all-zero block (they contribute
nothing to the sums below), so the model can score any (chemical, species,
duration) triplet whose chemical and species were each observed at least once
during training, but not entirely novel chemicals or species. For the dataset
used here $n=70{,}670$, $C=3{,}295$ chemicals, $1{,}267$ species, and
$p=4{,}879$ (namely $3{,}295+1{,}267+4+283+28+2$).

% =====================================================================
\subsection{The factorization-machine predictor}
\label{sec:fm}

The predictive mean for observation $i$ is the second-order factorization
machine
\begin{equation}
  \yhat_i(\theta) \;=\; w_0 \;+\; \sum_{j=1}^{p} w_j\,x_{ij}
  \;+\; \sum_{j=1}^{p}\sum_{j'>j}^{p} \langle \vv_j,\vv_{j'}\rangle\,x_{ij}x_{ij'},
  \label{eq:fm}
\end{equation}
with $\langle \vv_j,\vv_{j'}\rangle=\sum_{f=1}^{k} v_{jf}v_{j'f}$. The pairwise
interactions are not stored explicitly; they are represented in the low-rank
form of Eq.~\eqref{eq:fm}, which lets the model infer the interaction between a
chemical and a species that were never tested together. Using the standard FM
identity, the interaction sum is evaluated in $O(pk)$ time (in practice $O(k
N_{\mathrm{nz}})$ for a sparse $\xx_i$ with $N_{\mathrm{nz}}$ non-zeros) as
\begin{equation}
  \sum_{j<j'} \langle \vv_j,\vv_{j'}\rangle\,x_{ij}x_{ij'}
  \;=\; \tfrac{1}{2}\sum_{f=1}^{k}\!\left[
      \Big(\sum_{j=1}^{p} v_{jf}x_{ij}\Big)^{\!2}
      - \sum_{j=1}^{p} v_{jf}^{2}x_{ij}^{2}
  \right].
  \label{eq:fm-identity}
\end{equation}
The diagonal subtraction in Eq.~\eqref{eq:fm-identity} removes each feature's
self-interaction. A direct consequence, used repeatedly by the sampler, is that
$\yhat_i$ is an \emph{affine} function of each individual parameter when the
others are held fixed (Section~\ref{sec:sampler}).

% =====================================================================
\subsection{Full hierarchical generative model}
\label{sec:model}

The model is a Bayesian FM in which the single global noise precision of the
standard BFM is replaced by a per-chemical precision $\alpha_c$ under a shared
Gamma prior. The complete generative model is:

\paragraph{Likelihood.} Each observation is its FM mean plus Gaussian noise
whose precision is that of the observation's chemical,
\begin{equation}
  y_i \;\mid\; \xx_i,\theta,\alpha_{c_i}
  \;\sim\; \Norm\!\big(\yhat_i(\theta),\ \alpha_{c_i}^{-1}\big),
  \qquad i=1,\dots,n.
  \label{eq:like}
\end{equation}

\paragraph{Priors on the model parameters.} The global bias has a fixed Gaussian
prior; the weights and factors are drawn from shared populations whose mean and
precision are themselves learned:
\begin{align}
  w_0     &\sim \Norm(\mu_0,\ \gamma_0^{-1}), \\
  w_j     &\sim \Norm(\mu_w,\ \lambda_w^{-1}), & j&=1,\dots,p, \label{eq:prior-w}\\
  v_{jf}  &\sim \Norm(\mu_v,\ \lambda_v^{-1}), & j&=1,\dots,p,\ f=1,\dots,k. \label{eq:prior-v}
\end{align}

\paragraph{Hyperpriors on the populations.} There is one population for all
linear weights and one for all latent factors (as in the base BFM of
\citet{rendle2011bayesian}); each has a Gaussian hyperprior on its mean and a
Gamma hyperprior on its precision:
\begin{align}
  \mu_w &\sim \Norm(\mu_0,\ \gamma_0^{-1}), & \lambda_w &\sim \Gam(\alpha_\lambda,\ \beta_\lambda), \\
  \mu_v &\sim \Norm(\mu_0,\ \gamma_0^{-1}), & \lambda_v &\sim \Gam(\alpha_\lambda,\ \beta_\lambda).
\end{align}

\paragraph{Hierarchical noise model.} This is the methodological contribution.
Each chemical has its own noise precision, drawn from a Gamma prior whose rate is
shared across all chemicals and is itself given a Gamma hyperprior:
\begin{align}
  \alpha_c &\sim \Gam(a_0,\ b_0), & c&=1,\dots,C, \label{eq:prior-alpha}\\
  b_0      &\sim \Gam(c_0,\ d_0). \label{eq:prior-b0}
\end{align}
The shape $a_0$ is held fixed (set by empirical Bayes, Section~\ref{sec:eb});
the rate $b_0$ is resampled every sweep. Because $b_0$ is common to all
chemicals, it ties the per-chemical precisions together: a chemical with many
observations has its $\alpha_c$ set mainly by the scatter of its own residuals,
whereas a sparsely measured chemical is pooled towards the population-typical
noise level implied by $\Gam(a_0,b_0)$. This partial pooling is what keeps the
aleatoric estimate stable for data-poor chemicals.

The joint density factorises as
\begin{equation}
\begin{aligned}
  p\big(\theta,\{\alpha_c\},b_0,\mu_w,\lambda_w,\mu_v,\lambda_v \mid \mathbf{y}\big)
  \;\propto\;
  & \Big[\textstyle\prod_{i=1}^{n} \Norm\!\big(y_i \mid \yhat_i(\theta),\alpha_{c_i}^{-1}\big)\Big]
    \Norm(w_0\mid\mu_0,\gamma_0^{-1}) \\
  &\times \Big[\textstyle\prod_{j=1}^{p}\Norm(w_j\mid\mu_w,\lambda_w^{-1})\Big]
          \Big[\textstyle\prod_{j,f}\Norm(v_{jf}\mid\mu_v,\lambda_v^{-1})\Big] \\
  &\times \Norm(\mu_w\mid\mu_0,\gamma_0^{-1})\,\Gam(\lambda_w\mid\alpha_\lambda,\beta_\lambda)\,
          \Norm(\mu_v\mid\mu_0,\gamma_0^{-1})\,\Gam(\lambda_v\mid\alpha_\lambda,\beta_\lambda) \\
  &\times \Big[\textstyle\prod_{c=1}^{C}\Gam(\alpha_c\mid a_0,b_0)\Big]\,\Gam(b_0\mid c_0,d_0).
\end{aligned}
\label{eq:joint}
\end{equation}
This is the model drawn in plate notation in the main text: the fixed constants
$(\mu_0,\gamma_0,\alpha_\lambda,\beta_\lambda,a_0,c_0,d_0)$ set the hyperpriors;
$w_j$ and $v_{jf}$ share $(\mu_w,\lambda_w)$ and $(\mu_v,\lambda_v)$
respectively; the $\alpha_c$ share the rate $b_0$; and each $y_i$ is drawn from a
Gaussian centred at $\yhat_i$ with precision $\alpha_{c_i}$.

% =====================================================================
\subsection{Relationship to the standard Bayesian FM}
\label{sec:diff}

Our model departs from the BFM of \citet{rendle2011bayesian} in two respects.
\begin{enumerate}
  \item \textbf{Per-chemical noise (the substantive change).} The standard BFM
  assumes a single global precision $\alpha$ shared by all observations, with a
  Gamma prior and a Gamma-conjugate posterior. We replace it by $C$
  chemical-specific precisions $\alpha_c$ (Eq.~\ref{eq:prior-alpha}) under a
  shared, learned rate $b_0$ (Eq.~\ref{eq:prior-b0}). Equivalently, the global
  $\alpha$ in every conditional below becomes the observation-specific
  $\alpha_{c_i}$, and one extra conjugate draw for $b_0$ is added per sweep. This
  turns the noise term into a per-chemical, partially pooled quantity, which is
  the aleatoric component carried into the downstream analysis.

  \item \textbf{Independent Gaussian priors on the population means (a minor
  change).} The standard BFM couples each population mean to its precision
  through a Normal--Gamma prior, $\mu \mid \lambda \sim
  \Norm(\mu_0,(\gamma_0\lambda)^{-1})$, so the posterior precision of $\mu$
  scales with $\lambda$. We instead place an independent
  $\Norm(\mu_0,\gamma_0^{-1})$ prior on each mean (Eqs.~\ref{eq:prior-w}--\ref{eq:prior-v}),
  decoupling $\mu$ from $\lambda$; the mean does not then enter the update for
  $\lambda$, and the $\mu$-posterior precision is $\gamma_0+p\lambda$ rather than
  $(\gamma_0+p)\lambda$ (Section~\ref{sec:sampler}). With $p$ in the thousands
  this has negligible effect on inference, as \citet{rendle2011bayesian} also
  note for the hyperprior constants, but we state it for exactness.
\end{enumerate}
As in the reference, we use single-parameter (unblocked) Gibbs sampling, which
avoids the matrix inversions of blocked schemes and keeps each sweep linear in
$k$.

% =====================================================================
\subsection{Gibbs sampler and conditional posteriors}
\label{sec:sampler}

Posterior inference is by Gibbs sampling: each quantity is drawn in turn from its
full conditional given the current values of the others. All conditionals are
available in closed form because (i) the FM mean is affine in each single
parameter, giving Gaussian conditionals with Gaussian priors, and (ii) the
Gamma priors on the precisions are conjugate to the Gaussian likelihood.

Throughout we cache the running residual
\begin{equation}
  e_i \;=\; y_i - \yhat_i(\theta),
  \label{eq:resid}
\end{equation}
and the per-factor cache $q_{if}=\sum_{j=1}^{p} v_{jf}x_{ij}$, both updated by a
rank-one correction after each parameter is resampled. This is what makes a full
sweep $O(k N_{\mathrm{nz}})$, where $N_{\mathrm{nz}}$ is the number of non-zeros
in the design matrix. For a parameter with current value $\theta_\ell$ and
``slope'' $h_{i\ell}=\partial \yhat_i/\partial\theta_\ell$ (constant in
$\theta_\ell$ by multilinearity), the partial residual that adds back its current
contribution is $r_{i}=e_i+h_{i\ell}\theta_\ell$.

\subsubsection{Global bias $w_0$}
Here $h_{i0}=1$ for all $i$, so with prior $\Norm(\mu_0,\gamma_0^{-1})$,
\begin{equation}
  \sigma_{w_0}^2=\Big(\gamma_0+\sum_{i=1}^{n}\alpha_{c_i}\Big)^{-1},
  \qquad
  \mu_{w_0}=\sigma_{w_0}^2\Big(\gamma_0\mu_0+\sum_{i=1}^{n}\alpha_{c_i}\,(e_i+w_0)\Big),
  \qquad
  w_0 \sim \Norm(\mu_{w_0},\sigma_{w_0}^2).
\end{equation}
Then set $e_i \mathrel{+}= (w_0^{\mathrm{old}}-w_0^{\mathrm{new}})$ for all $i$.

\subsubsection{Weight population $(\mu_w,\lambda_w)$}
Treating the $p$ weights as i.i.d.\ $\Norm(\mu_w,\lambda_w^{-1})$ draws,
\begin{align}
  \mu_w &\sim \Norm\!\big(\mu_{\mu_w},\,\sigma_{\mu_w}^2\big),
  & \sigma_{\mu_w}^2 &=\big(\gamma_0+p\,\lambda_w\big)^{-1},
  & \mu_{\mu_w} &=\sigma_{\mu_w}^2\big(\gamma_0\mu_0+\lambda_w\textstyle\sum_{j=1}^{p}w_j\big), \\
  \lambda_w &\sim \Gam\!\Big(\alpha_\lambda+\tfrac{p}{2},\ \
     \beta_\lambda+\tfrac{1}{2}\textstyle\sum_{j=1}^{p}(w_j-\mu_w)^2\Big).
\end{align}

\subsubsection{Linear weights $w_j$}
For feature $j$, $h_{ij}=x_{ij}$, and only observations with $x_{ij}\neq 0$
contribute. With prior $\Norm(\mu_w,\lambda_w^{-1})$ and partial residual
$r_i=e_i+x_{ij}w_j$,
\begin{equation}
  \sigma_{w_j}^2=\Big(\lambda_w+\sum_{i:\,x_{ij}\neq0}\!\alpha_{c_i}\,x_{ij}^2\Big)^{-1},
  \quad
  \mu_{w_j}=\sigma_{w_j}^2\Big(\lambda_w\mu_w+\sum_{i:\,x_{ij}\neq0}\!\alpha_{c_i}\,x_{ij}\,r_i\Big),
  \quad
  w_j\sim\Norm(\mu_{w_j},\sigma_{w_j}^2).
\end{equation}
Then update $e_i \mathrel{+}= x_{ij}\,(w_j^{\mathrm{old}}-w_j^{\mathrm{new}})$.

\subsubsection{Factor population $(\mu_v,\lambda_v)$}
Treating all $pk$ factor entries as i.i.d.\ $\Norm(\mu_v,\lambda_v^{-1})$,
\begin{align}
  \mu_v &\sim \Norm\!\big(\mu_{\mu_v},\,\sigma_{\mu_v}^2\big),
  & \sigma_{\mu_v}^2 &=\big(\gamma_0+p\,k\,\lambda_v\big)^{-1},
  & \mu_{\mu_v} &=\sigma_{\mu_v}^2\big(\gamma_0\mu_0+\lambda_v\textstyle\sum_{j,f}v_{jf}\big), \\
  \lambda_v &\sim \Gam\!\Big(\alpha_\lambda+\tfrac{p k}{2},\ \
     \beta_\lambda+\tfrac{1}{2}\textstyle\sum_{j,f}(v_{jf}-\mu_v)^2\Big).
\end{align}

\subsubsection{Latent factors $v_{jf}$}
By the multilinearity of Eq.~\eqref{eq:fm-identity}, $\yhat_i$ is affine in
$v_{jf}$ with slope
\begin{equation}
  h_{ijf}=x_{ij}\big(q_{if}-x_{ij}v_{jf}\big),
  \qquad q_{if}=\sum_{j'=1}^{p} v_{j'f}x_{ij'},
\end{equation}
which does not depend on $v_{jf}$ (the $x_{ij}v_{jf}$ subtraction removes it).
With prior $\Norm(\mu_v,\lambda_v^{-1})$ and partial residual
$r_i=e_i+h_{ijf}v_{jf}$,
\begin{equation}
  \sigma_{v_{jf}}^2=\Big(\lambda_v+\sum_{i:\,x_{ij}\neq0}\!\alpha_{c_i}\,h_{ijf}^2\Big)^{-1},
  \quad
  \mu_{v_{jf}}=\sigma_{v_{jf}}^2\Big(\lambda_v\mu_v+\sum_{i:\,x_{ij}\neq0}\!\alpha_{c_i}\,h_{ijf}\,r_i\Big),
  \quad
  v_{jf}\sim\Norm(\mu_{v_{jf}},\sigma_{v_{jf}}^2).
\end{equation}
Then update both caches: $e_i \mathrel{+}= h_{ijf}\,(v_{jf}^{\mathrm{old}}-v_{jf}^{\mathrm{new}})$
and $q_{if}\mathrel{+}= x_{ij}\,(v_{jf}^{\mathrm{new}}-v_{jf}^{\mathrm{old}})$.

\subsubsection{Per-chemical noise precisions $\alpha_c$}
Let $N_c=|\{i:c_i=c\}|$ and $\mathrm{SSE}_c=\sum_{i:\,c_i=c} e_i^2$ be the
observation count and residual sum of squares of chemical $c$. Conjugacy of the
Gamma prior (Eq.~\ref{eq:prior-alpha}) to the Gaussian likelihood
(Eq.~\ref{eq:like}) gives
\begin{equation}
  \alpha_c \;\sim\; \Gam\!\Big(a_0+\tfrac{N_c}{2},\ \ b_0+\tfrac{1}{2}\,\mathrm{SSE}_c\Big),
  \qquad c=1,\dots,C.
  \label{eq:post-alpha}
\end{equation}
All $C$ precisions are drawn in one vectorised step using per-group sums of
$e_i^2$.

\subsubsection{Shared noise rate $b_0$}
The rate $b_0$ is conjugate to the $C$ Gamma-distributed precisions with common
shape $a_0$: combining $\prod_c \Gam(\alpha_c\mid a_0,b_0)\propto
b_0^{C a_0}e^{-b_0\sum_c\alpha_c}$ with the prior $\Gam(c_0,d_0)$ yields
\begin{equation}
  b_0 \;\sim\; \Gam\!\Big(c_0+C\,a_0,\ \ d_0+\textstyle\sum_{c=1}^{C}\alpha_c\Big).
  \label{eq:post-b0}
\end{equation}

\subsubsection{One sweep and sample collection}
A single Gibbs sweep applies, in order: $w_0$; $(\mu_w,\lambda_w)$ then every
$w_j$; $(\mu_v,\lambda_v)$ then every $v_{jf}$; the $\alpha_c$
(Eq.~\ref{eq:post-alpha}); and $b_0$ (Eq.~\ref{eq:post-b0}). The chain is run
for $n_{\mathrm{iter}}$ sweeps from an arbitrary initialisation
($w_0=0$, $w=\mathbf{0}$, $v_{jf}\sim\Norm(0,0.1^2)$, $\alpha_c=1$); the first
$n_{\mathrm{burn}}$ sweeps are discarded as burn-in and the remaining
$S=n_{\mathrm{iter}}-n_{\mathrm{burn}}$ sweeps are retained. Each retained sample
$s$ stores $\{w_0^{(s)},w^{(s)},v^{(s)},\{\alpha_c^{(s)}\},b_0^{(s)}\}$ and
corresponds to one complete, internally consistent parameterisation of the
model.

% =====================================================================
\subsection{Empirical-Bayes initialisation of the prior shape $a_0$}
\label{sec:eb}

The shared shape $a_0$ controls how strongly sparsely tested chemicals are
pooled towards the population noise level. Rather than fix it arbitrarily, we set
it by a method-of-moments empirical-Bayes pre-pass so that the prior matches the
dispersion of \emph{well-estimated} precisions:
\begin{enumerate}
  \item Run a short Gibbs pre-pass on the full training set with a non-informative,
  fixed noise prior ($a_0=1$, $b_0=1$, $b_0$ not learned) and collect the
  posterior draws $\{\alpha_c^{(s)}\}$.
  \item Restrict to data-rich chemicals ($N_c\ge N_{\min}$, with $N_{\min}=50$),
  whose posterior precisions are essentially data-driven, and form their
  posterior means $\hat{\alpha}_c=\tfrac{1}{S}\sum_s\alpha_c^{(s)}$.
  \item Fit a Gamma to $\{\hat{\alpha}_c\}$ by matching mean $m$ and variance
  $\nu$, and take the shape
  \begin{equation}
    a_0 \;=\; m^2/\nu.
  \end{equation}
  The corresponding rate is discarded, because the main run learns $b_0$ from its
  hyperprior (Eq.~\ref{eq:post-b0}).
\end{enumerate}
For the run reported in the paper this gives $a_0\approx1.75$; the main-run
posterior for the shared rate concentrates near $b_0\approx0.08$, so
$\Gam(a_0,b_0)$ is a weakly informative population prior on the per-chemical
precisions. The same EB $a_0$ is used for the cross-validation runs and for the
full-data model used to generate SSDs.

% =====================================================================
\subsection{Hyperparameter settings}
\label{sec:hyper}

Table~\ref{tab:hyper} lists the fixed constants and run settings. The mean- and
weight/factor-precision hyperpriors are the same trivial constants used by
\citet{rendle2011bayesian}; with $p$ in the thousands their influence on the
posterior is negligible. The noise-model hyperprior $\Gam(c_0,d_0)$ on $b_0$ and
the EB shape $a_0$ are specific to this work.

\begin{table}[htbp]
\centering
\caption{Hyperparameter and run settings.}
\label{tab:hyper}
\small
\begin{tabular}{lll}
\toprule
\textbf{Quantity} & \textbf{Symbol} & \textbf{Value} \\
\midrule
Factorization rank            & $k$                       & $32$ \\
Mean-hyperprior mean          & $\mu_0$                   & $0$ \\
Mean-hyperprior precision     & $\gamma_0$                & $1$ \\
Weight/factor precision prior & $\alpha_\lambda,\beta_\lambda$ & $1,\ 1$ \\
Noise-rate hyperprior         & $c_0,\ d_0$               & $1,\ 1$ \\
Noise-prior shape (empirical Bayes) & $a_0$               & $\approx 1.75$ (fixed after EB) \\
Noise-prior rate              & $b_0$                     & initialised at $1$, learned \\
Main run                      & $n_{\mathrm{iter}},\ n_{\mathrm{burn}}$ & $2000,\ 100$ \\
Retained samples              & $S$                       & $1900$ \\
EB pre-pass                   & $n_{\mathrm{iter}},\ n_{\mathrm{burn}}$ & $100,\ 50$ \\
EB data-rich threshold        & $N_{\min}$                & $50$ \\
Cross-validation              & folds                     & $5$ (grouped by triplet) \\
\bottomrule
\end{tabular}
\end{table}

% =====================================================================
\subsection{Posterior-predictive quantities}
\label{sec:predictive}

For a query $\xx^\ast$ belonging to chemical $c^\ast$, each retained sample gives
a deterministic FM prediction $\yhat^{(s)}(\xx^\ast)$ via Eq.~\eqref{eq:fm}. The
point prediction is the sample mean
\begin{equation}
  \yhat^\ast=\frac{1}{S}\sum_{s=1}^{S}\yhat^{(s)}(\xx^\ast).
\end{equation}
The predictive variance decomposes by the law of total variance,
$\Var(y^\ast\mid\mathcal{D})=\Var_\theta\!\big(\E[y^\ast\mid\theta]\big)
+\E_\theta\!\big[\Var(y^\ast\mid\theta)\big]$, whose two Monte-Carlo estimates are
the epistemic and aleatoric components:
\begin{align}
  \Var_{\mathrm{epi}}(\xx^\ast) &= \frac{1}{S}\sum_{s=1}^{S}\big(\yhat^{(s)}(\xx^\ast)-\yhat^\ast\big)^2,
  \label{eq:epi}\\
  \Var_{\mathrm{ale}}(c^\ast) &= \frac{1}{S}\sum_{s=1}^{S}\frac{1}{\alpha_{c^\ast}^{(s)}},
  \label{eq:ale}\\
  \Var_{\mathrm{tot}}(\xx^\ast,c^\ast) &= \Var_{\mathrm{epi}}(\xx^\ast)+\Var_{\mathrm{ale}}(c^\ast).
  \label{eq:tot}
\end{align}
The epistemic term is the spread of the mean across plausible fits and shrinks as
data constrain the parameters; the aleatoric term is the posterior-mean noise
variance of the chemical and does not vanish with more data. The aleatoric term
depends only on the chemical, so it is constant across species and durations for
a given chemical.

\paragraph{Calibration.} To assess whether $\Var_{\mathrm{tot}}$ is well
calibrated we form posterior-predictive draws that combine both components. For
each held-out observation and each sample $s$ we draw
\begin{equation}
  \tilde{y}^{(s)} \sim \Norm\!\big(\yhat^{(s)}(\xx^\ast),\ 1/\alpha_{c^\ast}^{(s)}\big),
\end{equation}
giving an $S$-point predictive sample. For a nominal central level $\rho$ we take
the empirical $[(1-\rho)/2,\ (1+\rho)/2]$ quantiles of these draws as the
credible interval, and report empirical coverage as the fraction of held-out
observations whose true value lies inside. Coverage is evaluated on the
out-of-fold observations (Section~\ref{sec:cv}), so predictions are scored by
models that did not train on them.

% =====================================================================
\subsection{SSD ensembles and hazardous concentrations}
\label{sec:ssd}

The Bayesian structure lets uncertainty propagate coherently into Species
Sensitivity Distributions. For a target chemical $c^\ast$ at a fixed duration,
and each retained sample $s$, we evaluate Eq.~\eqref{eq:fm} at $\theta^{(s)}$ for
all species $i=1,\dots,M$ and sort the resulting predictions ascending to obtain
one SSD curve $\big(\yhat^{(s)}_{(1)}\le\cdots\le\yhat^{(s)}_{(M)}\big)$ plotted
against the plotting positions $m/(M+1)$. Because every curve comes from a single
coherent parameter draw, the correlated epistemic structure across species is
preserved: the whole curve shifts and rescales as the shared parameters vary,
rather than species scattering independently. The ensemble of $S$ curves is
summarised pointwise by its median and central $95\%$ band.

A hazardous concentration $\mathrm{HC}x$ (the concentration protective of
$100-x\%$ of species) is read off each sampled curve at rank
$\lfloor x\,M/100\rfloor$, giving a posterior sample of $\mathrm{HC}x$ values;
its median and $2.5/97.5$ percentiles are the reported estimate and $95\%$
credible interval. Optionally, aleatoric noise can be injected per species by
adding a draw $\Norm\!\big(0,1/\alpha_{c^\ast}^{(s)}\big)$ to each prediction
before sorting; the SSD figures and HC intervals in the main text use the
epistemic ensemble (no added aleatoric noise) unless stated otherwise.

% =====================================================================
\subsection{Cross-validation and out-of-fold uncertainty}
\label{sec:cv}

Predictive accuracy and calibration are evaluated with grouped $5$-fold
cross-validation. Folds are formed by \emph{grouped} splitting on the
(chemical, species, duration) triplet identifier, so that all replicate
measurements of a held-out triplet are removed from training together and no
replicate of a scored triplet is ever seen during training. Each fold refits the
full model (including the hierarchical noise model, with the EB shape $a_0$ fixed
and $b_0$ learned) on its training partition and predicts the held-out partition;
concatenating the held-out predictions over folds yields the out-of-fold arrays
used for the RMSE, the uncertainty-versus-data analyses
(Eqs.~\ref{eq:epi}--\ref{eq:ale}), and the calibration curve. The full-data model
used to generate the SSDs and the all-chemical HC comparison is fit once on all
$n$ observations with the same settings.

Two practical points. First, the epistemic estimate of Eq.~\eqref{eq:epi} is only
meaningful for chemicals that appear in the training partition of a fold;
chemicals with fewer total observations than the number of folds can have folds
in which none of their rows are present, leaving their chemical-specific
parameters at the prior and deflating their apparent epistemic variance. Such
cold-start chemicals are excluded from the per-chemical
uncertainty-versus-data analyses in the main text. Second, the aleatoric
out-of-fold value for an observation is the sample mean of
$1/\alpha_{c}^{(s)}$ over that fold's retained samples, so it reflects the noise
precision learned from the fold's training data only.

% =====================================================================
\subsection{Computational complexity}
\label{sec:complexity}

Every conditional in Section~\ref{sec:sampler} is a single Gaussian or Gamma
draw, and the residual and factor caches make each parameter update a rank-one
correction touching only the non-zeros of the relevant feature column. One full
sweep therefore costs $O(k\,N_{\mathrm{nz}})$, linear in the factorization rank
and in the number of non-zero design-matrix entries, with an additional
$O(C)$ for the noise block (Eqs.~\ref{eq:post-alpha}--\ref{eq:post-b0}). This
matches the complexity of the unblocked BFM sampler of
\citet{rendle2011bayesian}; the per-chemical noise model adds only the vectorised
$\alpha_c$ draw and one scalar $b_0$ draw per sweep, so the hierarchical
extension does not change the asymptotic cost.

.
% Relies on the preamble of main.tex (amsmath, booktabs, etc. are already
% loaded there; amssymb is added there for \Real below). Uses the shared
% references.bib already cited by the rest of the paper.
% =====================================================================

% ---- convenience macros ----
\newcommand{\yhat}{\hat{y}}
\newcommand{\Norm}{\mathcal{N}}
\newcommand{\Gam}{\mathrm{Gamma}}
\newcommand{\Real}{\mathbb{R}}
\newcommand{\E}{\mathbb{E}}
\newcommand{\Var}{\mathrm{Var}}
\newcommand{\vv}{\mathbf{v}}
\newcommand{\xx}{\mathbf{x}}

\section{Supporting Information}
\label{sec:si}

This Supporting Information gives the full mathematical specification of the
Hierarchical Bayesian Factorization Machine (BFM) summarised in the main text.
It is written for a statistical or machine-learning audience and documents the
model exactly as implemented: the predictor, the complete hierarchical
generative model, the Gibbs sampler and every conditional posterior, the
empirical-Bayes initialisation, the hyperparameter settings, and the
posterior-predictive quantities used for calibration, Species Sensitivity
Distributions (SSDs), and hazardous concentrations. The predictive core is the
factorization machine (FM) of \citet{rendle2010fm} cast in the Bayesian form of
\citet{rendle2011bayesian}; Section~\ref{sec:diff} states precisely how our
model differs from theirs. Symbols follow the main text and are collected in
Table~\ref{tab:notation}.

% =====================================================================
\subsection{Notation}
\label{sec:notation}

\begin{table}[htbp]
\centering
\caption{Notation used throughout the Supporting Information.}
\label{tab:notation}
\small
\begin{tabular}{ll}
\toprule
\textbf{Symbol} & \textbf{Meaning} \\
\midrule
$n$ & number of observations (rows) \\
$p$ & number of features (design-matrix columns) \\
$k$ & number of latent factors per feature (factorization rank) \\
$C$ & number of chemicals (noise groups) \\
$\xx_i \in \Real^{p}$ & feature vector of observation $i$ \\
$x_{ij}$ & $j$-th entry of $\xx_i$ \\
$y_i$ & centred log-transformed LC50 of observation $i$ \\
$c_i \in \{1,\dots,C\}$ & chemical (group) index of observation $i$ \\
$w_0$ & global bias \\
$w_j$ & linear weight of feature $j$ \\
$v_{jf}$ & $f$-th latent factor of feature $j$; $\vv_j=(v_{j1},\dots,v_{jk})$ \\
$\alpha_c$ & noise precision of chemical $c$ (aleatoric term) \\
$\mu_w,\lambda_w$ & mean and precision of the shared prior on the weights $w_j$ \\
$\mu_v,\lambda_v$ & mean and precision of the shared prior on the factors $v_{jf}$ \\
$a_0,b_0$ & shape and rate of the shared prior on the $\alpha_c$ \\
$\mu_0,\gamma_0$ & mean and precision of the Gaussian mean-hyperpriors \\
$\alpha_\lambda,\beta_\lambda$ & shape and rate of the Gamma prior on $\lambda_w,\lambda_v$ \\
$c_0,d_0$ & shape and rate of the Gamma hyperprior on $b_0$ \\
$\theta$ & all model parameters $\{w_0,\{w_j\},\{v_{jf}\}\}$ \\
$S$ & number of retained (post-burn-in) posterior samples \\
$\yhat^{(s)}$ & prediction from posterior sample $s$ \\
\bottomrule
\end{tabular}
\end{table}

We write $\Norm(\mu,\sigma^2)$ for a Gaussian with mean $\mu$ and variance
$\sigma^2$, and $\Gam(a,b)$ for a Gamma distribution with shape $a$ and rate $b$
(so that its mean is $a/b$). All precisions are inverse variances.

% =====================================================================
\subsection{Data representation and feature encoding}
\label{sec:features}

Each observation is a single laboratory LC50 measurement for a
(chemical, species, duration) combination. The target $y_i$ is the
log-transformed concentration in mg/L, centred by subtracting the global mean
$\bar{y}$ of the training set; the mean is added back to recover log-mg/L units
when reporting predictions. Repeated measurements of the same triplet are kept
as separate rows rather than averaged, so that the per-chemical noise model
(Section~\ref{sec:model}) sees the within-triplet spread directly.

The feature vector $\xx_i$ concatenates five one-hot categorical blocks and two
standardised numerical features:
\begin{itemize}
  \item \textbf{chemical} identity (one-hot over $C$ chemicals),
  \item \textbf{species} identity (one-hot),
  \item \textbf{test duration} (one-hot over $\{24,48,72,96\}$\,h),
  \item \textbf{taxonomic family} and \textbf{taxonomic class} of the species (two one-hot blocks),
  \item \textbf{numerical}: log-transformed molecular weight and the RDKit
        octanol--water partition coefficient (\texttt{clogp}), median-imputed
        and standardised to zero mean and unit variance.
\end{itemize}
Unseen categories at prediction time map to an all-zero block (they contribute
nothing to the sums below), so the model can score any (chemical, species,
duration) triplet whose chemical and species were each observed at least once
during training, but not entirely novel chemicals or species. For the dataset
used here $n=70{,}670$, $C=3{,}295$ chemicals, $1{,}267$ species, and
$p=4{,}879$ (namely $3{,}295+1{,}267+4+283+28+2$).

% =====================================================================
\subsection{The factorization-machine predictor}
\label{sec:fm}

The predictive mean for observation $i$ is the second-order factorization
machine
\begin{equation}
  \yhat_i(\theta) \;=\; w_0 \;+\; \sum_{j=1}^{p} w_j\,x_{ij}
  \;+\; \sum_{j=1}^{p}\sum_{j'>j}^{p} \langle \vv_j,\vv_{j'}\rangle\,x_{ij}x_{ij'},
  \label{eq:fm}
\end{equation}
with $\langle \vv_j,\vv_{j'}\rangle=\sum_{f=1}^{k} v_{jf}v_{j'f}$. The pairwise
interactions are not stored explicitly; they are represented in the low-rank
form of Eq.~\eqref{eq:fm}, which lets the model infer the interaction between a
chemical and a species that were never tested together. Using the standard FM
identity, the interaction sum is evaluated in $O(pk)$ time (in practice $O(k
N_{\mathrm{nz}})$ for a sparse $\xx_i$ with $N_{\mathrm{nz}}$ non-zeros) as
\begin{equation}
  \sum_{j<j'} \langle \vv_j,\vv_{j'}\rangle\,x_{ij}x_{ij'}
  \;=\; \tfrac{1}{2}\sum_{f=1}^{k}\!\left[
      \Big(\sum_{j=1}^{p} v_{jf}x_{ij}\Big)^{\!2}
      - \sum_{j=1}^{p} v_{jf}^{2}x_{ij}^{2}
  \right].
  \label{eq:fm-identity}
\end{equation}
The diagonal subtraction in Eq.~\eqref{eq:fm-identity} removes each feature's
self-interaction. A direct consequence, used repeatedly by the sampler, is that
$\yhat_i$ is an \emph{affine} function of each individual parameter when the
others are held fixed (Section~\ref{sec:sampler}).

% =====================================================================
\subsection{Full hierarchical generative model}
\label{sec:model}

The model is a Bayesian FM in which the single global noise precision of the
standard BFM is replaced by a per-chemical precision $\alpha_c$ under a shared
Gamma prior. The complete generative model is:

\paragraph{Likelihood.} Each observation is its FM mean plus Gaussian noise
whose precision is that of the observation's chemical,
\begin{equation}
  y_i \;\mid\; \xx_i,\theta,\alpha_{c_i}
  \;\sim\; \Norm\!\big(\yhat_i(\theta),\ \alpha_{c_i}^{-1}\big),
  \qquad i=1,\dots,n.
  \label{eq:like}
\end{equation}

\paragraph{Priors on the model parameters.} The global bias has a fixed Gaussian
prior; the weights and factors are drawn from shared populations whose mean and
precision are themselves learned:
\begin{align}
  w_0     &\sim \Norm(\mu_0,\ \gamma_0^{-1}), \\
  w_j     &\sim \Norm(\mu_w,\ \lambda_w^{-1}), & j&=1,\dots,p, \label{eq:prior-w}\\
  v_{jf}  &\sim \Norm(\mu_v,\ \lambda_v^{-1}), & j&=1,\dots,p,\ f=1,\dots,k. \label{eq:prior-v}
\end{align}

\paragraph{Hyperpriors on the populations.} There is one population for all
linear weights and one for all latent factors (as in the base BFM of
\citet{rendle2011bayesian}); each has a Gaussian hyperprior on its mean and a
Gamma hyperprior on its precision:
\begin{align}
  \mu_w &\sim \Norm(\mu_0,\ \gamma_0^{-1}), & \lambda_w &\sim \Gam(\alpha_\lambda,\ \beta_\lambda), \\
  \mu_v &\sim \Norm(\mu_0,\ \gamma_0^{-1}), & \lambda_v &\sim \Gam(\alpha_\lambda,\ \beta_\lambda).
\end{align}

\paragraph{Hierarchical noise model.} This is the methodological contribution.
Each chemical has its own noise precision, drawn from a Gamma prior whose rate is
shared across all chemicals and is itself given a Gamma hyperprior:
\begin{align}
  \alpha_c &\sim \Gam(a_0,\ b_0), & c&=1,\dots,C, \label{eq:prior-alpha}\\
  b_0      &\sim \Gam(c_0,\ d_0). \label{eq:prior-b0}
\end{align}
The shape $a_0$ is held fixed (set by empirical Bayes, Section~\ref{sec:eb});
the rate $b_0$ is resampled every sweep. Because $b_0$ is common to all
chemicals, it ties the per-chemical precisions together: a chemical with many
observations has its $\alpha_c$ set mainly by the scatter of its own residuals,
whereas a sparsely measured chemical is pooled towards the population-typical
noise level implied by $\Gam(a_0,b_0)$. This partial pooling is what keeps the
aleatoric estimate stable for data-poor chemicals.

The joint density factorises as
\begin{equation}
\begin{aligned}
  p\big(\theta,\{\alpha_c\},b_0,\mu_w,\lambda_w,\mu_v,\lambda_v \mid \mathbf{y}\big)
  \;\propto\;
  & \Big[\textstyle\prod_{i=1}^{n} \Norm\!\big(y_i \mid \yhat_i(\theta),\alpha_{c_i}^{-1}\big)\Big]
    \Norm(w_0\mid\mu_0,\gamma_0^{-1}) \\
  &\times \Big[\textstyle\prod_{j=1}^{p}\Norm(w_j\mid\mu_w,\lambda_w^{-1})\Big]
          \Big[\textstyle\prod_{j,f}\Norm(v_{jf}\mid\mu_v,\lambda_v^{-1})\Big] \\
  &\times \Norm(\mu_w\mid\mu_0,\gamma_0^{-1})\,\Gam(\lambda_w\mid\alpha_\lambda,\beta_\lambda)\,
          \Norm(\mu_v\mid\mu_0,\gamma_0^{-1})\,\Gam(\lambda_v\mid\alpha_\lambda,\beta_\lambda) \\
  &\times \Big[\textstyle\prod_{c=1}^{C}\Gam(\alpha_c\mid a_0,b_0)\Big]\,\Gam(b_0\mid c_0,d_0).
\end{aligned}
\label{eq:joint}
\end{equation}
This is the model drawn in plate notation in the main text: the fixed constants
$(\mu_0,\gamma_0,\alpha_\lambda,\beta_\lambda,a_0,c_0,d_0)$ set the hyperpriors;
$w_j$ and $v_{jf}$ share $(\mu_w,\lambda_w)$ and $(\mu_v,\lambda_v)$
respectively; the $\alpha_c$ share the rate $b_0$; and each $y_i$ is drawn from a
Gaussian centred at $\yhat_i$ with precision $\alpha_{c_i}$.

% =====================================================================
\subsection{Relationship to the standard Bayesian FM}
\label{sec:diff}

Our model departs from the BFM of \citet{rendle2011bayesian} in two respects.
\begin{enumerate}
  \item \textbf{Per-chemical noise (the substantive change).} The standard BFM
  assumes a single global precision $\alpha$ shared by all observations, with a
  Gamma prior and a Gamma-conjugate posterior. We replace it by $C$
  chemical-specific precisions $\alpha_c$ (Eq.~\ref{eq:prior-alpha}) under a
  shared, learned rate $b_0$ (Eq.~\ref{eq:prior-b0}). Equivalently, the global
  $\alpha$ in every conditional below becomes the observation-specific
  $\alpha_{c_i}$, and one extra conjugate draw for $b_0$ is added per sweep. This
  turns the noise term into a per-chemical, partially pooled quantity, which is
  the aleatoric component carried into the downstream analysis.

  \item \textbf{Independent Gaussian priors on the population means (a minor
  change).} The standard BFM couples each population mean to its precision
  through a Normal--Gamma prior, $\mu \mid \lambda \sim
  \Norm(\mu_0,(\gamma_0\lambda)^{-1})$, so the posterior precision of $\mu$
  scales with $\lambda$. We instead place an independent
  $\Norm(\mu_0,\gamma_0^{-1})$ prior on each mean (Eqs.~\ref{eq:prior-w}--\ref{eq:prior-v}),
  decoupling $\mu$ from $\lambda$; the mean does not then enter the update for
  $\lambda$, and the $\mu$-posterior precision is $\gamma_0+p\lambda$ rather than
  $(\gamma_0+p)\lambda$ (Section~\ref{sec:sampler}). With $p$ in the thousands
  this has negligible effect on inference, as \citet{rendle2011bayesian} also
  note for the hyperprior constants, but we state it for exactness.
\end{enumerate}
As in the reference, we use single-parameter (unblocked) Gibbs sampling, which
avoids the matrix inversions of blocked schemes and keeps each sweep linear in
$k$.

% =====================================================================
\subsection{Gibbs sampler and conditional posteriors}
\label{sec:sampler}

Posterior inference is by Gibbs sampling: each quantity is drawn in turn from its
full conditional given the current values of the others. All conditionals are
available in closed form because (i) the FM mean is affine in each single
parameter, giving Gaussian conditionals with Gaussian priors, and (ii) the
Gamma priors on the precisions are conjugate to the Gaussian likelihood.

Throughout we cache the running residual
\begin{equation}
  e_i \;=\; y_i - \yhat_i(\theta),
  \label{eq:resid}
\end{equation}
and the per-factor cache $q_{if}=\sum_{j=1}^{p} v_{jf}x_{ij}$, both updated by a
rank-one correction after each parameter is resampled. This is what makes a full
sweep $O(k N_{\mathrm{nz}})$, where $N_{\mathrm{nz}}$ is the number of non-zeros
in the design matrix. For a parameter with current value $\theta_\ell$ and
``slope'' $h_{i\ell}=\partial \yhat_i/\partial\theta_\ell$ (constant in
$\theta_\ell$ by multilinearity), the partial residual that adds back its current
contribution is $r_{i}=e_i+h_{i\ell}\theta_\ell$.

\subsubsection{Global bias $w_0$}
Here $h_{i0}=1$ for all $i$, so with prior $\Norm(\mu_0,\gamma_0^{-1})$,
\begin{equation}
  \sigma_{w_0}^2=\Big(\gamma_0+\sum_{i=1}^{n}\alpha_{c_i}\Big)^{-1},
  \qquad
  \mu_{w_0}=\sigma_{w_0}^2\Big(\gamma_0\mu_0+\sum_{i=1}^{n}\alpha_{c_i}\,(e_i+w_0)\Big),
  \qquad
  w_0 \sim \Norm(\mu_{w_0},\sigma_{w_0}^2).
\end{equation}
Then set $e_i \mathrel{+}= (w_0^{\mathrm{old}}-w_0^{\mathrm{new}})$ for all $i$.

\subsubsection{Weight population $(\mu_w,\lambda_w)$}
Treating the $p$ weights as i.i.d.\ $\Norm(\mu_w,\lambda_w^{-1})$ draws,
\begin{align}
  \mu_w &\sim \Norm\!\big(\mu_{\mu_w},\,\sigma_{\mu_w}^2\big),
  & \sigma_{\mu_w}^2 &=\big(\gamma_0+p\,\lambda_w\big)^{-1},
  & \mu_{\mu_w} &=\sigma_{\mu_w}^2\big(\gamma_0\mu_0+\lambda_w\textstyle\sum_{j=1}^{p}w_j\big), \\
  \lambda_w &\sim \Gam\!\Big(\alpha_\lambda+\tfrac{p}{2},\ \
     \beta_\lambda+\tfrac{1}{2}\textstyle\sum_{j=1}^{p}(w_j-\mu_w)^2\Big).
\end{align}

\subsubsection{Linear weights $w_j$}
For feature $j$, $h_{ij}=x_{ij}$, and only observations with $x_{ij}\neq 0$
contribute. With prior $\Norm(\mu_w,\lambda_w^{-1})$ and partial residual
$r_i=e_i+x_{ij}w_j$,
\begin{equation}
  \sigma_{w_j}^2=\Big(\lambda_w+\sum_{i:\,x_{ij}\neq0}\!\alpha_{c_i}\,x_{ij}^2\Big)^{-1},
  \quad
  \mu_{w_j}=\sigma_{w_j}^2\Big(\lambda_w\mu_w+\sum_{i:\,x_{ij}\neq0}\!\alpha_{c_i}\,x_{ij}\,r_i\Big),
  \quad
  w_j\sim\Norm(\mu_{w_j},\sigma_{w_j}^2).
\end{equation}
Then update $e_i \mathrel{+}= x_{ij}\,(w_j^{\mathrm{old}}-w_j^{\mathrm{new}})$.

\subsubsection{Factor population $(\mu_v,\lambda_v)$}
Treating all $pk$ factor entries as i.i.d.\ $\Norm(\mu_v,\lambda_v^{-1})$,
\begin{align}
  \mu_v &\sim \Norm\!\big(\mu_{\mu_v},\,\sigma_{\mu_v}^2\big),
  & \sigma_{\mu_v}^2 &=\big(\gamma_0+p\,k\,\lambda_v\big)^{-1},
  & \mu_{\mu_v} &=\sigma_{\mu_v}^2\big(\gamma_0\mu_0+\lambda_v\textstyle\sum_{j,f}v_{jf}\big), \\
  \lambda_v &\sim \Gam\!\Big(\alpha_\lambda+\tfrac{p k}{2},\ \
     \beta_\lambda+\tfrac{1}{2}\textstyle\sum_{j,f}(v_{jf}-\mu_v)^2\Big).
\end{align}

\subsubsection{Latent factors $v_{jf}$}
By the multilinearity of Eq.~\eqref{eq:fm-identity}, $\yhat_i$ is affine in
$v_{jf}$ with slope
\begin{equation}
  h_{ijf}=x_{ij}\big(q_{if}-x_{ij}v_{jf}\big),
  \qquad q_{if}=\sum_{j'=1}^{p} v_{j'f}x_{ij'},
\end{equation}
which does not depend on $v_{jf}$ (the $x_{ij}v_{jf}$ subtraction removes it).
With prior $\Norm(\mu_v,\lambda_v^{-1})$ and partial residual
$r_i=e_i+h_{ijf}v_{jf}$,
\begin{equation}
  \sigma_{v_{jf}}^2=\Big(\lambda_v+\sum_{i:\,x_{ij}\neq0}\!\alpha_{c_i}\,h_{ijf}^2\Big)^{-1},
  \quad
  \mu_{v_{jf}}=\sigma_{v_{jf}}^2\Big(\lambda_v\mu_v+\sum_{i:\,x_{ij}\neq0}\!\alpha_{c_i}\,h_{ijf}\,r_i\Big),
  \quad
  v_{jf}\sim\Norm(\mu_{v_{jf}},\sigma_{v_{jf}}^2).
\end{equation}
Then update both caches: $e_i \mathrel{+}= h_{ijf}\,(v_{jf}^{\mathrm{old}}-v_{jf}^{\mathrm{new}})$
and $q_{if}\mathrel{+}= x_{ij}\,(v_{jf}^{\mathrm{new}}-v_{jf}^{\mathrm{old}})$.

\subsubsection{Per-chemical noise precisions $\alpha_c$}
Let $N_c=|\{i:c_i=c\}|$ and $\mathrm{SSE}_c=\sum_{i:\,c_i=c} e_i^2$ be the
observation count and residual sum of squares of chemical $c$. Conjugacy of the
Gamma prior (Eq.~\ref{eq:prior-alpha}) to the Gaussian likelihood
(Eq.~\ref{eq:like}) gives
\begin{equation}
  \alpha_c \;\sim\; \Gam\!\Big(a_0+\tfrac{N_c}{2},\ \ b_0+\tfrac{1}{2}\,\mathrm{SSE}_c\Big),
  \qquad c=1,\dots,C.
  \label{eq:post-alpha}
\end{equation}
All $C$ precisions are drawn in one vectorised step using per-group sums of
$e_i^2$.

\subsubsection{Shared noise rate $b_0$}
The rate $b_0$ is conjugate to the $C$ Gamma-distributed precisions with common
shape $a_0$: combining $\prod_c \Gam(\alpha_c\mid a_0,b_0)\propto
b_0^{C a_0}e^{-b_0\sum_c\alpha_c}$ with the prior $\Gam(c_0,d_0)$ yields
\begin{equation}
  b_0 \;\sim\; \Gam\!\Big(c_0+C\,a_0,\ \ d_0+\textstyle\sum_{c=1}^{C}\alpha_c\Big).
  \label{eq:post-b0}
\end{equation}

\subsubsection{One sweep and sample collection}
A single Gibbs sweep applies, in order: $w_0$; $(\mu_w,\lambda_w)$ then every
$w_j$; $(\mu_v,\lambda_v)$ then every $v_{jf}$; the $\alpha_c$
(Eq.~\ref{eq:post-alpha}); and $b_0$ (Eq.~\ref{eq:post-b0}). The chain is run
for $n_{\mathrm{iter}}$ sweeps from an arbitrary initialisation
($w_0=0$, $w=\mathbf{0}$, $v_{jf}\sim\Norm(0,0.1^2)$, $\alpha_c=1$); the first
$n_{\mathrm{burn}}$ sweeps are discarded as burn-in and the remaining
$S=n_{\mathrm{iter}}-n_{\mathrm{burn}}$ sweeps are retained. Each retained sample
$s$ stores $\{w_0^{(s)},w^{(s)},v^{(s)},\{\alpha_c^{(s)}\},b_0^{(s)}\}$ and
corresponds to one complete, internally consistent parameterisation of the
model.

% =====================================================================
\subsection{Empirical-Bayes initialisation of the prior shape $a_0$}
\label{sec:eb}

The shared shape $a_0$ controls how strongly sparsely tested chemicals are
pooled towards the population noise level. Rather than fix it arbitrarily, we set
it by a method-of-moments empirical-Bayes pre-pass so that the prior matches the
dispersion of \emph{well-estimated} precisions:
\begin{enumerate}
  \item Run a short Gibbs pre-pass on the full training set with a non-informative,
  fixed noise prior ($a_0=1$, $b_0=1$, $b_0$ not learned) and collect the
  posterior draws $\{\alpha_c^{(s)}\}$.
  \item Restrict to data-rich chemicals ($N_c\ge N_{\min}$, with $N_{\min}=50$),
  whose posterior precisions are essentially data-driven, and form their
  posterior means $\hat{\alpha}_c=\tfrac{1}{S}\sum_s\alpha_c^{(s)}$.
  \item Fit a Gamma to $\{\hat{\alpha}_c\}$ by matching mean $m$ and variance
  $\nu$, and take the shape
  \begin{equation}
    a_0 \;=\; m^2/\nu.
  \end{equation}
  The corresponding rate is discarded, because the main run learns $b_0$ from its
  hyperprior (Eq.~\ref{eq:post-b0}).
\end{enumerate}
For the run reported in the paper this gives $a_0\approx1.75$; the main-run
posterior for the shared rate concentrates near $b_0\approx0.08$, so
$\Gam(a_0,b_0)$ is a weakly informative population prior on the per-chemical
precisions. The same EB $a_0$ is used for the cross-validation runs and for the
full-data model used to generate SSDs.

% =====================================================================
\subsection{Hyperparameter settings}
\label{sec:hyper}

Table~\ref{tab:hyper} lists the fixed constants and run settings. The mean- and
weight/factor-precision hyperpriors are the same trivial constants used by
\citet{rendle2011bayesian}; with $p$ in the thousands their influence on the
posterior is negligible. The noise-model hyperprior $\Gam(c_0,d_0)$ on $b_0$ and
the EB shape $a_0$ are specific to this work.

\begin{table}[htbp]
\centering
\caption{Hyperparameter and run settings.}
\label{tab:hyper}
\small
\begin{tabular}{lll}
\toprule
\textbf{Quantity} & \textbf{Symbol} & \textbf{Value} \\
\midrule
Factorization rank            & $k$                       & $32$ \\
Mean-hyperprior mean          & $\mu_0$                   & $0$ \\
Mean-hyperprior precision     & $\gamma_0$                & $1$ \\
Weight/factor precision prior & $\alpha_\lambda,\beta_\lambda$ & $1,\ 1$ \\
Noise-rate hyperprior         & $c_0,\ d_0$               & $1,\ 1$ \\
Noise-prior shape (empirical Bayes) & $a_0$               & $\approx 1.75$ (fixed after EB) \\
Noise-prior rate              & $b_0$                     & initialised at $1$, learned \\
Main run                      & $n_{\mathrm{iter}},\ n_{\mathrm{burn}}$ & $2000,\ 100$ \\
Retained samples              & $S$                       & $1900$ \\
EB pre-pass                   & $n_{\mathrm{iter}},\ n_{\mathrm{burn}}$ & $100,\ 50$ \\
EB data-rich threshold        & $N_{\min}$                & $50$ \\
Cross-validation              & folds                     & $5$ (grouped by triplet) \\
\bottomrule
\end{tabular}
\end{table}

% =====================================================================
\subsection{Posterior-predictive quantities}
\label{sec:predictive}

For a query $\xx^\ast$ belonging to chemical $c^\ast$, each retained sample gives
a deterministic FM prediction $\yhat^{(s)}(\xx^\ast)$ via Eq.~\eqref{eq:fm}. The
point prediction is the sample mean
\begin{equation}
  \yhat^\ast=\frac{1}{S}\sum_{s=1}^{S}\yhat^{(s)}(\xx^\ast).
\end{equation}
The predictive variance decomposes by the law of total variance,
$\Var(y^\ast\mid\mathcal{D})=\Var_\theta\!\big(\E[y^\ast\mid\theta]\big)
+\E_\theta\!\big[\Var(y^\ast\mid\theta)\big]$, whose two Monte-Carlo estimates are
the epistemic and aleatoric components:
\begin{align}
  \Var_{\mathrm{epi}}(\xx^\ast) &= \frac{1}{S}\sum_{s=1}^{S}\big(\yhat^{(s)}(\xx^\ast)-\yhat^\ast\big)^2,
  \label{eq:epi}\\
  \Var_{\mathrm{ale}}(c^\ast) &= \frac{1}{S}\sum_{s=1}^{S}\frac{1}{\alpha_{c^\ast}^{(s)}},
  \label{eq:ale}\\
  \Var_{\mathrm{tot}}(\xx^\ast,c^\ast) &= \Var_{\mathrm{epi}}(\xx^\ast)+\Var_{\mathrm{ale}}(c^\ast).
  \label{eq:tot}
\end{align}
The epistemic term is the spread of the mean across plausible fits and shrinks as
data constrain the parameters; the aleatoric term is the posterior-mean noise
variance of the chemical and does not vanish with more data. The aleatoric term
depends only on the chemical, so it is constant across species and durations for
a given chemical.

\paragraph{Calibration.} To assess whether $\Var_{\mathrm{tot}}$ is well
calibrated we form posterior-predictive draws that combine both components. For
each held-out observation and each sample $s$ we draw
\begin{equation}
  \tilde{y}^{(s)} \sim \Norm\!\big(\yhat^{(s)}(\xx^\ast),\ 1/\alpha_{c^\ast}^{(s)}\big),
\end{equation}
giving an $S$-point predictive sample. For a nominal central level $\rho$ we take
the empirical $[(1-\rho)/2,\ (1+\rho)/2]$ quantiles of these draws as the
credible interval, and report empirical coverage as the fraction of held-out
observations whose true value lies inside. Coverage is evaluated on the
out-of-fold observations (Section~\ref{sec:cv}), so predictions are scored by
models that did not train on them.

% =====================================================================
\subsection{SSD ensembles and hazardous concentrations}
\label{sec:ssd}

The Bayesian structure lets uncertainty propagate coherently into Species
Sensitivity Distributions. For a target chemical $c^\ast$ at a fixed duration,
and each retained sample $s$, we evaluate Eq.~\eqref{eq:fm} at $\theta^{(s)}$ for
all species $i=1,\dots,M$ and sort the resulting predictions ascending to obtain
one SSD curve $\big(\yhat^{(s)}_{(1)}\le\cdots\le\yhat^{(s)}_{(M)}\big)$ plotted
against the plotting positions $m/(M+1)$. Because every curve comes from a single
coherent parameter draw, the correlated epistemic structure across species is
preserved: the whole curve shifts and rescales as the shared parameters vary,
rather than species scattering independently. The ensemble of $S$ curves is
summarised pointwise by its median and central $95\%$ band.

A hazardous concentration $\mathrm{HC}x$ (the concentration protective of
$100-x\%$ of species) is read off each sampled curve at rank
$\lfloor x\,M/100\rfloor$, giving a posterior sample of $\mathrm{HC}x$ values;
its median and $2.5/97.5$ percentiles are the reported estimate and $95\%$
credible interval. Optionally, aleatoric noise can be injected per species by
adding a draw $\Norm\!\big(0,1/\alpha_{c^\ast}^{(s)}\big)$ to each prediction
before sorting; the SSD figures and HC intervals in the main text use the
epistemic ensemble (no added aleatoric noise) unless stated otherwise.

% =====================================================================
\subsection{Cross-validation and out-of-fold uncertainty}
\label{sec:cv}

Predictive accuracy and calibration are evaluated with grouped $5$-fold
cross-validation. Folds are formed by \emph{grouped} splitting on the
(chemical, species, duration) triplet identifier, so that all replicate
measurements of a held-out triplet are removed from training together and no
replicate of a scored triplet is ever seen during training. Each fold refits the
full model (including the hierarchical noise model, with the EB shape $a_0$ fixed
and $b_0$ learned) on its training partition and predicts the held-out partition;
concatenating the held-out predictions over folds yields the out-of-fold arrays
used for the RMSE, the uncertainty-versus-data analyses
(Eqs.~\ref{eq:epi}--\ref{eq:ale}), and the calibration curve. The full-data model
used to generate the SSDs and the all-chemical HC comparison is fit once on all
$n$ observations with the same settings.

Two practical points. First, the epistemic estimate of Eq.~\eqref{eq:epi} is only
meaningful for chemicals that appear in the training partition of a fold;
chemicals with fewer total observations than the number of folds can have folds
in which none of their rows are present, leaving their chemical-specific
parameters at the prior and deflating their apparent epistemic variance. Such
cold-start chemicals are excluded from the per-chemical
uncertainty-versus-data analyses in the main text. Second, the aleatoric
out-of-fold value for an observation is the sample mean of
$1/\alpha_{c}^{(s)}$ over that fold's retained samples, so it reflects the noise
precision learned from the fold's training data only.

% =====================================================================
\subsection{Computational complexity}
\label{sec:complexity}

Every conditional in Section~\ref{sec:sampler} is a single Gaussian or Gamma
draw, and the residual and factor caches make each parameter update a rank-one
correction touching only the non-zeros of the relevant feature column. One full
sweep therefore costs $O(k\,N_{\mathrm{nz}})$, linear in the factorization rank
and in the number of non-zero design-matrix entries, with an additional
$O(C)$ for the noise block (Eqs.~\ref{eq:post-alpha}--\ref{eq:post-b0}). This
matches the complexity of the unblocked BFM sampler of
\citet{rendle2011bayesian}; the per-chemical noise model adds only the vectorised
$\alpha_c$ draw and one scalar $b_0$ draw per sweep, so the hierarchical
extension does not change the asymptotic cost.

